% ============================================================================
% 【已退役正文／RETIRED MANUSCRIPT】
% 本文件属于由英国算例／ALNS 旧稿迁移、继而采用 China81 与 MV-HGS-SP 的旧 V2 时代完整正文。
% 退役原因：项目现已采用现行 GCI-DMM-VRP 模型与一条串行改进链；
% 本文件的旧算法身份、模型口径、比较对象和实验结果均不再代表当前项目。
% 现行正文：docs/paper_gci_dmm_vrp_20260804/paper_main_foundation_v2_20260812.md
% 任何人不得引用、迁移或改写本文件中的实验数字作为当前论文证据。
% 本文件仅为历史追溯保留；除本声明外，原正文及其数字未作改动。
% ============================================================================

% ================================================================
% ReSETP 论文 V2 (2026-07-20 全面重写版)
% 结构模板: 逐段对齐陈雨蝶 2025 (系统工程理论与实践)
% 基座文献: 邱莹莹 2024 硕士学位论文 (混合车队+动态需求)
% 实验宇宙: China81 (京津冀/珠三角/成渝, 人民币, 48时段分时电价与碳强度)
% 算法: HGS-FRVCP (成熟混合遗传搜索底座 + frvcpy 固定路线充电求解 + 机制动作提名)
% 数学来源: %% MATH_REF: <来源论文>
% ================================================================

\documentclass{setp-new}
\usepackage{fontspec}
\setmainfont{Times New Roman}
\setsansfont{Arial}

% --- CJK 字体设置 (保留原稿配置) ---
\IfFontExistsTF{Songti SC}{\newcommand{\setpSongFont}{Songti SC}}{\newcommand{\setpSongFont}{SimSun}}
\newif\ifsetpLocalNotoHei
\IfFileExists{/Users/zhouleixishu/Library/Fonts/NotoSansCJKsc-Regular.otf}{%
  \setpLocalNotoHeitrue}{%
  \setpLocalNotoHeifalse}
\ifsetpLocalNotoHei\else
\IfFontExistsTF{SimHei}{\newcommand{\setpHeiFont}{SimHei}}{%
  \IfFontExistsTF{Noto Sans CJK SC}{\newcommand{\setpHeiFont}{Noto Sans CJK SC}}{%
    \PackageError{ReSETP}{No approved Simplified Chinese Heiti font found}{Install SimHei or Noto Sans CJK SC before compiling}%
    \newcommand{\setpHeiFont}{Noto Sans CJK SC}}}
\fi
\IfFontExistsTF{Songti SC}{\newcommand{\setpFangSongFont}{Songti SC}}{\newcommand{\setpFangSongFont}{FangSong}}
\ifsetpLocalNotoHei
\setCJKmainfont{\setpSongFont}
\setCJKsansfont[Path=/Users/zhouleixishu/Library/Fonts/]{NotoSansCJKsc-Regular.otf}
\setCJKmonofont{\setpSongFont}
\setCJKfamilyfont{zhsong}{\setpSongFont}
\setCJKfamilyfont{zhhei}[Path=/Users/zhouleixishu/Library/Fonts/]{NotoSansCJKsc-Regular.otf}
\setCJKfamilyfont{zhfs}{\setpFangSongFont}
\else
\setCJKmainfont[BoldFont=\setpHeiFont]{\setpSongFont}
\setCJKsansfont{\setpHeiFont}
\setCJKmonofont{\setpSongFont}
\setCJKfamilyfont{zhsong}[BoldFont=\setpHeiFont]{\setpSongFont}
\setCJKfamilyfont{zhhei}{\setpHeiFont}
\setCJKfamilyfont{zhfs}{\setpFangSongFont}
\fi
\providecommand{\songti}{}
\providecommand{\heiti}{}
\providecommand{\fangsong}{}
\renewcommand{\songti}{\CJKfamily{zhsong}}
\renewcommand{\heiti}{\CJKfamily{zhhei}}
\renewcommand{\fangsong}{\CJKfamily{zhfs}}
\XeTeXlinebreaklocale "zh"
\XeTeXlinebreakskip = 0pt plus 1pt

% --- 数学与排版宏包 ---
\usepackage{mathtools}
\usepackage{graphicx}
\usepackage{booktabs}
\usepackage{array}
\usepackage{tabularx}
\usepackage{longtable}
\usepackage{supertabular}
\usepackage{makecell}
\usepackage{multirow}
\usepackage{hhline}
\usepackage{calc}
\usepackage{ragged2e}
\usepackage{float}
\usepackage{needspace}
\usepackage{url}
\usepackage{tikz}
\usetikzlibrary{arrows.meta,backgrounds,fit,positioning,shapes.geometric}

% --- 表格样式 (保留原稿) ---
\newcommand{\setpsymtabsetup}{\fontsize{9pt}{10.5pt}\selectfont\renewcommand{\arraystretch}{0.95}\setlength{\tabcolsep}{2.0pt}\setlength{\aboverulesep}{0.4pt}\setlength{\belowrulesep}{0.4pt}\hbadness=10000}
\newcommand{\setptabsetup}{\fontsize{9pt}{10.5pt}\selectfont\renewcommand{\arraystretch}{0.95}\setlength{\tabcolsep}{2.2pt}\setlength{\aboverulesep}{0.4pt}\setlength{\belowrulesep}{0.4pt}\hbadness=10000}

% --- 紧凑数学间距 ---
\newcommand{\setpcompactmathspacing}{%
  \setlength{\abovedisplayskip}{3pt plus .6pt minus .6pt}%
  \setlength{\belowdisplayskip}{3pt plus .6pt minus .6pt}%
  \setlength{\abovedisplayshortskip}{2pt plus .5pt minus .5pt}%
  \setlength{\belowdisplayshortskip}{2pt plus .5pt minus .5pt}%
  \setlength{\jot}{1.5pt}%
  \setlength{\arraycolsep}{2.2pt}%
}

% --- 列类型 ---
\newcolumntype{S}{>{\centering\arraybackslash}m{0.115\linewidth}}
\newcolumntype{T}{>{\centering\arraybackslash}m{0.365\linewidth}}
\newcolumntype{Y}{>{\Centering\arraybackslash}p{0.105\linewidth}}
\newcolumntype{J}{>{\Centering\arraybackslash}p{0.375\linewidth}}
\newcolumntype{U}{>{\Centering\arraybackslash}p{0.088\linewidth}}
\newcolumntype{V}{>{\RaggedRight\arraybackslash}p{0.186\linewidth}}
\newcolumntype{Z}[1]{>{\centering\arraybackslash}m{#1}}

% --- 文档初始化 ---
\AtBeginDocument{%
  \raggedbottom
  \setpcompactmathspacing
  \setlength{\floatsep}{4pt plus .8pt minus .8pt}%
  \setlength{\textfloatsep}{5pt plus .8pt minus 1pt}%
  \setlength{\intextsep}{4pt plus .8pt minus .8pt}%
  \setlength{\LTpre}{1pt}%
  \setlength{\LTpost}{1pt}%
  \setlength{\heavyrulewidth}{0.4pt}%
  \setlength{\lightrulewidth}{0.3pt}%
  \setlength{\cmidrulewidth}{0.3pt}%
}
\AtBeginEnvironment{equation}{\setpcompactmathspacing}
\AtBeginEnvironment{gather}{\setpcompactmathspacing}
\AtBeginEnvironment{align}{\setpcompactmathspacing}
\AtBeginEnvironment{alignat}{\setpcompactmathspacing}

\renewcommand{\figurename}{图}
\renewcommand{\tablename}{表}
\newcounter{none}
\newcommand{\pandocbounded}[1]{#1}
\providecommand{\tightlist}{\setlength{\itemsep}{0pt}\setlength{\parskip}{0pt}}
\DeclareRobustCommand{\TBD}[1]{%
  \begingroup\setlength{\fboxsep}{1.2pt}%
  \colorbox{black!12}{\textcolor{black!65}{\textsf{\footnotesize[待填：#1]}}}%
  \endgroup}
\setlength{\parindent}{2em}
\setlength{\emergencystretch}{3em}
\setlength{\headheight}{32pt}
\raggedbottom
\allowdisplaybreaks

% ================================================================
% 论文元数据
% ================================================================
\jolname{系统工程理论与实践}{Systems Engineering --- Theory \& Practice}
\Volume{2026}{5}{0}{0}
\PageNum{1}
\PaperDOI{}
\Category{U492.3}{A}

% --- 中文信息 ---
\TitleMark{周雷习书等: 时变碳强度下多车场动态协同路径建模与算法}
\Title{时变碳强度下多车场动态协同路径建模与算法}
\Author{周雷习书, 干宏程, 邱莹莹}{上海理工大学 管理学院, 上海 200093}

% --- 中文摘要 (陈雨蝶模板: 背景→挑战→模型→算法→实验维度→价值) ---
\Abstract{%
配送车队电动化并不自动带来减排: 电动车的排放转移至充电环节,
而分时电价与电网碳强度在时间上并不同向, 仅按电费最省安排充电反而会推高充电排放.
针对这一问题, 构建时变电网碳强度下多家企业协作的动态混合车队多趟路径优化模型,
以可行配送趟—实体车排班两层结构刻画多趟衔接、趟间按需补电、非线性充电、
成员结算参与和滚动状态继承, 并将运营成本与碳成本纳入同一核算口径.
求解上直接采用成熟混合遗传搜索作为搜索底座,
在其上构造统一完整评价、实体车多趟补全、碳感知充电排程与动态状态继承等问题特化组件.
数值实验沿混合车队、真实补能、时变碳强度、多车场协同、收益公平与动态运行的顺序逐层展开,
逐项检验各现实因素是否改变最优配送组织与路径选择.
结果显示, 打开车场归属边界的协同重组使降本与减碳同向,
而充电时刻与车型指派两层的降本与减碳方向相反, 需由碳价对齐, 且两层的对齐水平并不相同;
相关结论可为配送企业安排补能与车型、为政策制定者设定碳价区间提供量化依据.
}

\Keywords{车辆路径问题; 多车场协同配送; 混合车队; 时变电网碳强度; 动态需求; 非线性充电; 混合遗传搜索}

% --- 英文信息 ---
\ETitle{Multi-depot collaborative routing model and algorithm under time-varying carbon intensity}
\EAuthor{ZHOU Leixishu, GAN Hongcheng, QIU Yingying}{Business School, University of Shanghai for Science and Technology, Shanghai 200093, China}
\EAbstract{%
Regional and intercity delivery is complicated by independent depot operations, mixed fleets, charging periods misaligned with grid carbon intensity, and dynamic orders.
This study formulates a collaborative multi-depot mixed-fleet routing model with dynamic demand and time-varying charging emissions.
A two-layer structure links feasible delivery trips to physical-vehicle schedules.
It represents multi-trip connections, inter-trip charging, settlement-based participation, and rolling-state inheritance while accounting for operating and carbon costs.
A mature hybrid genetic search is adopted as the search kernel, on which a Hybrid Genetic Search with a Fixed-Route Vehicle Charging Problem solver (HGS-FRVCP) is built.
Its problem-specific components comprise full-solution evaluation, physical-vehicle multi-trip completion, carbon-aware charging scheduling, and dynamic state inheritance.
Numerical experiments proceed in the order of mixed fleet, realistic recharging, time-varying carbon intensity, multi-depot collaboration, surplus fairness, and dynamic operation.
Each factor is examined for whether it alters the optimal delivery organization and routing rather than only the reported cost and emissions.
The results quantify trade-offs relevant to depot collaboration, charging-time selection, and carbon-policy design.
}
\EKeywords{vehicle routing problem; collaborative multi-depot delivery; mixed fleet;
time-varying grid carbon intensity; dynamic demand; nonlinear charging; hybrid genetic search}

\begin{document}
\maketitle
\clearpage

% ================================================================
% SECTION 1: 引言 (按陈雨蝶8段式)
% ================================================================

\section{引言}

%% --- 段1: 政策宣言 ---
2020年, 中国领导人在第七十五届联合国大会上承诺"2030年碳达峰,2060年碳中和", 简称"双碳目标".
国际能源署数据显示, 2023年我国碳排放总量约126亿吨, 交通运输占比约10\%.
交通运输减排因而成为实现双碳目标的重要环节.
《2030年前碳达峰行动方案》进一步提出推动运输装备低碳转型和配送集约化\cite{ref:55},
将协同降碳和集约配送上升为国家行动目标.

%% --- 段2: 行业痛点+学者研究 ---
长期以来, 城市与城际配送的主力车型是燃油车.
传统配送路径优化主要解决客户归属、车辆调度与路线排序,
目标是在按时完成配送的前提下尽量减少行驶距离与运营成本.
在减碳约束下, 这一目标结构不再充分: 车辆若始终燃油驱动,
仅在既有燃油车路径上继续压缩距离, 可获得的减排空间有限.
配送系统因此转向电动化.
但受购置成本、补能条件与车辆残值的限制, 车队不可能一次性完成替换,
燃油车与电动车将在相当长时期内共存, 企业实际面对的是混合车队.
由此产生第一个决策问题: 两类车型的续航、能耗与成本结构均不同,
何种配送任务应交由电动车、何种应保留给燃油车,
以及车型构成改变之后, 原有的最优路径是否仍然最优.
这些车场又往往分属不同的配送企业, 各企业拥有各自的车场、车队与客户责任,
降本减排空间既来自企业内部的车型选择与补能安排, 也来自企业之间的协作配送.
现有研究主要从充电调度、路径选择和车辆配置三个方面降低运营成本与排放.
在分时充电减排方面, Li等\cite{ref:52}基于上海3777辆电动车11个月的运行数据,
在不改变出行计划的条件下量化有序充电的减排潜力.
Hu等\cite{ref:53}将动态电价、节点碳流和碳交易用于电动车群充电调度.
Zhang等\cite{ref:54}联合优化行驶路径与充电站选择, 以缩短出行时间和平衡充电站负荷为目标.
Cheng等\cite{ref:50}在车辆可用时间、站端容量和电池约束下直接优化充电碳排.

%% --- 段3: VRP理论视角过渡 ---
电动车一旦真正投入运行, 后续问题依次显现:
真实补能过程并非恒功率, 荷电状态越高充入越慢, 补能时长会反过来决定一条路线能否按时完成;
补能过程刻画真实之后, 又须回答用电本身是否低碳——同样一度电在不同时段充入, 对应的电网碳排并不相同;
而何时可充又取决于何时发车、服务哪些客户、何时返场与在何处补能.
再往外, 限制效率的可能已不是车辆与补能, 而是车场之间人为划定的客户归属边界;
边界一旦打开, 系统总收益的改善又未必被每个参与方分享.
上述问题可统一表述为车辆路径问题(Vehicle Routing Problem, VRP)的扩展形式:
混合车队VRP、多车场VRP和动态VRP分别刻画车型差异、跨场资源配置与订单变化,
三者的耦合构成区域—城际配送的主要决策结构.

%% --- 段4: 子领域1 —— 混合车队低碳路径 ---
在混合车队低碳路径方面, Goeke和Schneider\cite{ref:27}建立了油电混合车队路径模型,
利用包含速度、坡度和载重的实际能耗函数替代线性距离能耗假设,
并设计自适应大邻域搜索算法求解.
李得成等\cite{ref:6}基于分支定价算法进一步研究车型配置与路径的联合决策.
陈婉茹等\cite{ref:23}将碳交易引入多中心混合车队的路径与速度决策.
Keskin和\c{C}atay\cite{ref:45}通过部分充电策略扩大电动车路径的可行决策空间.
Diefenbach等\cite{ref:46}表明非线性充电曲线会进一步改变充电与车辆排程.
Kullman等\cite{ref:73}将固定路线充电子问题封装为可嵌入更大路径搜索的开源精确求解工具frvcpy.
这些研究为车辆能耗和碳成本核算提供了基础,
但大多同时改变路线、车辆和充电安排,
难以单独判断减排是来自配送结构变化, 还是来自既定排班内的充电时刻调整.

%% --- 段5: 子领域2 —— 多车场协同 ---
在多车场协同方面, Wang等\cite{ref:7}考察了多车场电动车共享充电设施的路径协同.
Wang等\cite{ref:25}进一步在时变路网下研究多车场协同与资源共享.
Fern\'{a}ndez等\cite{ref:49}直接将共享客户作为协同路径的核心决策.
系统总收益增加并不表示每个参与方都获益.
Wang等\cite{ref:7}在协同路径优化后采用Shapley值分配合作收益,
Soriano等\cite{ref:8}则在多车场路径中直接考虑运营利润公平.
Shi等\cite{ref:9}进一步联合讨论协同运输计划与利润分配,
Agussurja等\cite{ref:10}研究了维持稳定且公平分配所需的最小补贴.
客户原责任与地理服务边界的关系还会影响合作空间:
责任划分越接近地理边界, 独立经营可能越有效,
留给跨场重组的增量空间反而越小.

%% --- 段6: 子领域3 —— 动态车辆路径 ---
在动态车辆路径方面, Psaraftis\cite{ref:14}较早系统讨论了动态信息下的路径决策.
李阳等\cite{ref:18}利用周期性重优化处理动态需求.
Gendreau等\cite{ref:22}采用并行禁忌搜索进行实时路径与调度.
邱莹莹\cite{ref:71}在低碳背景下研究考虑动态需求的混合车队配送路径,
采用定时与定量双触发的批处理策略处理新增、取消和改量订单.
在多车场多趟场景中, Zhen等\cite{ref:44}将配送趟生成和实体车衔接分别处理.
姜广田等\cite{ref:51}在绿色物流配送下研究多车型动态车辆路径优化.
但是, 多车场混合车队的滚动重规划不能只更新客户集合,
还必须继承车辆状态、已执行任务及已发生的成本和排放.
若缺少这些状态, 重优化方案可能撤回已执行路径、重复使用实体车或重复结算历史排放.

%% --- 段7: 综合+研究不足 ---
相关研究还将碳交易机制、时变路网和动态调度纳入车辆路径模型.
陈雨蝶等\cite{ref:48}在双碳背景下将多中心、时变路网和多类配送约束纳入同一冷链物流模型,
表明综合机制建模可以拓展问题的现实边界, 同时也会提高求解复杂度.
1) 混合车队与协同配送研究多把电动车视为零排放或按固定排放因子折算,
电力供给的时段差异没有进入决策.
例如Shi等\cite{ref:93}在混合车队、协同配送与分时电价的联合模型中,
碳排放目标只对燃油车行驶弧求和, 并据此得出全部改用电动车即可使碳排放降为零的结论;
在该设定下充电时刻只改变电费, 不改变排放.
然而分时电价与电网碳强度在时间上并不同步:
本文所用中国代表电网日数据显示, 谷段八小时的平均碳强度为0.618\,kgCO$_2$e/kWh,
高于平段的0.334与峰段的0.367, 二者相关系数为$-0.563$;
仅按电费最省安排充电, 反而会把充电推向全天碳强度最高的时段.
充电时刻因而不应只作为成本变量, 而应作为进入路径决策的减排杠杆.
2) 有序充电研究多假定出行计划不变,
对配送时间窗、实体车多趟和动态事件共同形成的可移动充电窗口讨论不足;
配送路径研究又常将路线与充电联合改变, 两类作用不易分离.
3) 协同研究分别关注系统成本、运营利润或事后收益分配,
联盟总收益、成员参与和子联盟稳定性仍需在同一价值函数下衔接.
4) 降本与减碳分属配送企业与政策制定者两类主体的关切,
二者在充电时刻、车型指派与配送结构上并不同向;
既有研究多以单一目标或加权目标处理, 对何种碳价水平才能使二者对齐缺少可操作结论.

%% --- 段8: 本文工作+创新点 ---
针对上述缺口, 本文以降本与减碳为双重目标, 面向配送企业与政策制定者两类主体,
建立时变电网碳强度下多家企业协作的动态混合车队多趟路径优化模型,
并在成熟混合遗传搜索之上设计面向该问题的特化求解组件对其求解.
论文主要创新点如下:
1) 同时考虑多车场协同、混合车队、实体车多趟、分时充电排放、成员结算参与和动态订单等诸多因素,
建立可行配送趟—实体车排班两层模型,
用趟间按需补电和滚动状态继承衔接静态计划与动态执行,
调整相应参数后, 该模型可退化为单车场、单趟或纯燃油车队等经典VRP形式.
2) 将48时段电网碳强度和分时电价引入路径、车型、充电站与充电时刻的联合决策,
并以固定配送排班下的充电择时作为补充检验,
分别量化联合优化的排放--成本权衡和既定排班内充电时刻的边际减排作用.
3) 在成熟混合遗传搜索之上设计面向该问题的特化求解组件,
包括完整方案评价、实体车多趟补全、碳感知充电排程和动态状态继承;
候选方案的可行性与成本均按完整模型核算, 不以简化路线代价近似替代.
4) 围绕时变电网碳强度、混合车队和动态需求组织数值实验,
并以多车场协同配送、成员收益分配、实体车多趟和非线性充电作为支撑结构,
分别考察充电时刻、车型指派和动态发车三个决策层对运营成本与碳排放的影响.

为进一步明确本文与相关研究的要素边界, 比较结果见表\ref{tab:literature-comparison}.

% --- 文献比较表 ---
\begin{table}[H]
\centering
\caption{动态协同混合车队相关文献比较}
\label{tab:literature-comparison}
\setptabsetup
\begin{tabular*}{\linewidth}{@{\extracolsep{\fill}}lcccccc@{}}
\toprule
文献 & 多车场 & 混合车队 & 实体车多趟 & 动态事件 & 分时充电排放 & 成员收益条件\\
\midrule
Goeke和Schneider\cite{ref:27} & --- & $\checkmark$ & --- & --- & --- & ---\\
Zhen等\cite{ref:44} & $\checkmark$ & --- & $\checkmark$ & --- & --- & ---\\
陈婉茹等\cite{ref:23} & $\checkmark$ & $\checkmark$ & --- & --- & --- & ---\\
Wang等\cite{ref:7} & $\checkmark$ & --- & --- & --- & --- & ---\\
Soriano等\cite{ref:8} & $\checkmark$ & --- & --- & --- & --- & $\checkmark$\\
邱莹莹\cite{ref:71} & --- & $\checkmark$ & --- & $\checkmark$ & --- & ---\\
姜广田等\cite{ref:51} & --- & --- & --- & $\checkmark$ & --- & ---\\
Li等\cite{ref:52} & --- & --- & --- & --- & $\checkmark$ & ---\\
本文 & $\checkmark$ & $\checkmark$ & $\checkmark$ & $\checkmark$ & $\checkmark$ & $\checkmark$\\
\bottomrule
\end{tabular*}
\end{table}

% ================================================================
% SECTION 2: 模型建立 (按陈雨蝶4子节式)
% ================================================================

\section{模型建立}

\subsection{问题描述}

本文研究多家配送企业协作配送时的混合车队路径、充电和成员结算问题,
并在动态订单到达后滚动更新尚未执行的任务.
各企业独立拥有车场、车队和客户责任, 协作意味着允许跨企业调用车辆并重新分工,
因而所得方案能否被各方接受, 须在结算层单独判定.

本文模型具体表述为:
多家配送企业各自拥有一个车场及其车队, 并各自负责一部分客户;
协作配送允许各企业共享客户和车辆资源, 联合为全部客户提供服务;
车队由燃油车(CV)和电动车(EV)混合组成, 每辆车只属于一个企业的车场, 但允许服务其他企业的客户;
车辆可在一个运营日内连续执行多趟任务,
电动车途中可访问共享充电站, 也可在首趟出发前和相邻趟之间在车场按需补电;
实体车队规模有限且每辆车按日计一次启用成本, 因而一辆车通常需承担多趟任务,
其日间返场装货的间隔正是车场补能的可用窗口, 补能机会由此与配送排班直接绑定;
电网碳强度和电价均按48个半小时时段给定,
充电发生的时段不同, 充电电费和间接碳排放也随之不同;
订单新增、取消或改量后, 系统只调整尚未执行的任务, 已完成或已发车的路径前缀不可撤回.
模型的优化目标是在满足载重、时间窗、电量和站点容量等运行约束的前提下,
以运营成本和碳结算成本之和最小化确定车辆路径、车型选择和充电安排;
合作成员的参与性在方案生成后的结算层判定.

模型假设条件如下:
1) 在每个规划时刻, 已到达订单的客户、车场、车辆信息已知, 新订单在运营过程中动态出现.
2) 每个客户由一辆车服务一次, 需求不可拆分.
3) 车辆在配送全程不超过所属车型的最大载重.
4) 车辆行驶速度设为常值, 能耗按综合模式排放模型(CMEM)系能耗函数\cite{ref:27,ref:36}计算.
5) 客户具有硬时间窗, 车辆可提前到达后等待, 但不得迟于时间窗上界.
6) 充电功率为电池荷电状态的分段常函数(恒功率为其退化情形),
每次充电为一个连续区间, 允许部分充电, 时段电量由充电曲线和电网时段的重叠确定.
7) 时段电网碳强度和分时电价为外生参数, 企业为价格与碳强度的接受者.
8) 客户责任按编号顺序均分给各配送企业, 该归属为公开披露的构造情景, 分法参照饶卫振等\cite{ref:91};
运营成本先按实际执行车辆所属企业归集;
协作路线优化完成后, 再依据各子联盟的独立经营成本分配协作收益.
核仁给出唯一且联盟稳定的结算结果, Shapley值作为成员贡献的比较基准.
模型符号说明如表\ref{tab:symbols}所示.

\begin{table}[H]
\centering
\caption{符号说明}
\label{tab:symbols}
\setpsymtabsetup
\begin{tabular*}{\linewidth}{@{\extracolsep{\fill}}cJcJ@{}}
\toprule
符号 & 说明 & 符号 & 说明\\
\midrule
$D$ & 车场集合 & $c_D$ & 空气阻力系数\\
$N$ & 客户集合 $N=\{1,2,\dots,n\}$, 其中$n$是客户数目 & $c_R$ & 滚动阻力系数\\
$S$ & 共享充电站集合 & $\rho^a$ & 空气密度(kg/m$^3$)\\
$T$ & 电网时段集合 $T=\{1,2,\dots,48\}$ & $g_0$ & 重力加速度(m/s$^2$)\\
$K$ & 实体车辆集合 & $\alpha_k^e$ & 车辆$k$的电耗修正系数\\
$K_d$ & 车场$d$的车辆集合 & $\beta_{0k}^g,\beta_{1k}^g$ & 车辆$k$的油耗系数\\
$V_r$ & 趟$r$的节点集合 & $B$ & 电池容量(kWh)\\
$V_r^c$ & 趟$r$的客户子集 $V_r^c=V_r\cap N$ & $B^{\min}$ & 电池最低允许电量(kWh)\\
$\Omega_k$ & 车辆$k$的可行排班模式集合 & $\lambda^g$ & 燃油排放因子(kgCO$_2$e/L)\\
$\mathcal R_{kp}$ & 模式$p$的配送趟集合 & $\pi_s$ & 站点$s$的额定充电功率(kW)\\
$\mathcal Q_{kp}$ & 模式$p$的充电会话集合 & $C_s$ & 站点$s$的充电桩数\\
$d_r^+,d_r^-$ & 趟$r$的起讫车场副本 & $b_l,\kappa_l$ & 充电曲线第$l$段SOC断点与相对功率\\
$[\underline T_t,\overline T_t]$ & 时段$t$的起止时刻(s) & $\widehat\gamma_t,\gamma_t$ & 时段$t$的预测与结算碳强度(kgCO$_2$e/kWh)\\
$q_i$ & 客户$i$的需求量(kg) & $p^{\mathrm{car}}$ & 单位碳价格(元/kg)\\
$Q_k$ & 车辆$k$的最大载重(kg) & $CE$ & 碳配额(kg)\\
$[\underline t_i,\overline t_i]$ & 客户$i$的时间窗 & $c^{\mathrm{fix}}$ & 每趟派遣固定成本(元/趟)\\
$\sigma_i$ & 客户$i$的服务时长(s) & $c_k^{\mathrm{km}}$ & 车辆$k$的单位里程成本(元/km)\\
$d_{ij}$ & 节点$i$与节点$j$之间的距离(m) & $p^f$ & 单位燃油价格(元/L)\\
$t_{ij}$ & 弧$(i,j)$的行驶时间(s) & $p_{s,t}^{e}$ & 站点$s$在时段$t$的电价(元/kWh)\\
$v$ & 车辆行驶速度(m/s) & $c^{\mathrm{occ}}$ & 公共站单位占用时间成本(元/min)\\
$m_k$ & 车辆$k$的整备质量(kg) & $c^{\mathrm{tr}}$ & 单位跨场协调成本(元/客户)\\
$A_k^f$ & 车辆$k$的迎风面积(m$^2$) & $\rho$ & 单位货量收入(元/kg)\\
$d(k)$ & 车辆$k$所属车场 & $v(A)$ & 企业联盟$A$的可转移收益(元)\\
$d_i^0$ & 客户$i$的原责任车场 & $\Pi_d^0$ & 车场$d$独立经营的基准收益(元)\\
$C(A),R(A)$ & 联盟$A$的配送成本与收入(元)
& $\phi_d$ & 车场$d$的Shapley收益分配(元)\\
$x_d^{\nu}$ & 车场$d$的核仁收益分配(元)
& $e(A,x)$ & 联盟$A$对分配$x$的不满意额(元)\\
$n_{kp}$ & 模式$p$的配送趟数 & $\overline W$ & 动态触发的累计需求上限(kg)\\
$D_{kp}$ & 模式$p$的总里程(km) & $T^{\mathrm{int}}$ & 动态触发的时间间隔(s)\\
$F_{kp}$ & 模式$p$的总油耗(L) & $M$ & 足够大的正数\\
$H_{kp}^{S},\widehat E_{kp},E_{kp}$ & 模式$p$的公共站占用时长(min)、预测与结算碳排放(kgCO$_2$e) & $L_{ir}$ & 离开节点$i$时的载重(kg)\\
$O_{skpt}$ & 模式$p$在站$s$、时段$t$的在充会话数 & $T_{ir}$ & 节点$i$的服务开始时刻(s)\\
$A_{ikp}$ & 0-1变量, 如果模式$p$服务客户$i$, 则为1, 否则为0 & $\varepsilon_{ir}$ & 到达节点$i$时的电量(kWh)\\
$x_{ijr}$ & 0-1变量, 如果趟$r$使用弧$(i,j)$, 则为1, 否则为0 & $Y_{ir}$ & 节点$i$处的补电量(kWh)\\
$a_{ir}$ & 0-1变量, 如果趟$r$服务客户$i$, 则为1, 否则为0 & $a_q^{\mathrm{ch}},\Delta_q$ & 会话$q$的开始时刻与时长(s)\\
$z_{kp}$ & 0-1变量, 如果车辆$k$选择模式$p$, 则为1, 否则为0 & $E_q^{\mathrm{ch}},y_{qt}$ & 会话$q$的补电量与其在时段$t$的补电量(kWh)\\
\bottomrule
\end{tabular*}
\end{table}

\subsection{车辆能耗模型}

为了准确刻画载重和车型对能耗的影响,
本文参照混合车队路径优化研究\cite{ref:27}采用CMEM系能耗函数\cite{ref:36,ref:47},
车辆在弧$(i,j)$上以给定速度$v$匀速行驶,
牵引功率由空气阻力、滚动阻力和载重共同决定:
\begin{equation}
P_{ij}=\frac12c_D\rho^a A_k^f v^3+(m_k+u_{ij})g_0c_R v.
\label{eq:power}
\end{equation}
其中$u_{ij}$为弧段载重(kg), 由载重变量$L_{ir}$按弧确定.
式(\ref{eq:power})中$P_{ij}$的单位为W. 为使电量与电池状态变量的单位一致,
电动车弧段电耗按$1\ \mathrm{kWh}=3.6\times10^{6}\ \mathrm{J}$折算为kWh;
燃油车弧段油耗由排放系数直接给出体积单位L. 二者分别为
\begin{align}
e_{ij}&=\frac{\alpha_k^e}{3.6\times10^{6}}P_{ij}\frac{d_{ij}}{v}, \label{eq:electricity}\\
f_{ij}&=(\beta_{0k}^g+\beta_{1k}^gP_{ij})\frac{d_{ij}}{v}. \label{eq:fuel}
\end{align}
其中$\alpha_k^e$为无量纲的电驱动系统效率系数.
在常值速度下, 式(\ref{eq:electricity})退化为关于载重线性的按里程电耗式:
\begin{equation}
e_{ij}=(g_v+\omega^E u_{ij})d_{ij},\quad
g_v=\frac{\alpha_k^e}{3.6\times10^{6}}\Bigl(\frac12c_D\rho^a A_k^f v^2+m_k g_0 c_R\Bigr),\quad
\omega^E=\frac{\alpha_k^e g_0 c_R}{3.6\times10^{6}}.
\label{eq:linear_energy}
\end{equation}
其中$d_{ij}$以m计, $g_v$与$\omega^E$分别为kWh/m和kWh/(kg$\cdot$m),
因而$e_{ij}$与$\varepsilon_{ir}$、$Y_{ir}$、$B$同为kWh, 式(\ref{eq:batt_prop})的电量传播在量纲上闭合.
该线性形式使电量传播约束保持线性, 燃油油耗同理退化.

\subsection{充电与分时碳排放模型}

传统电动车路径研究常将充电近似为恒功率满充,
这与现实中动力电池高荷电状态下功率下降、车辆按需部分补电的运行方式不符.
鉴于此, 本文参照部分充电与非线性充电研究\cite{ref:45,ref:46},
将站点$s$的充电功率表示为电池荷电状态(SOC)的分段常函数.
设$L$个SOC分段$[b_{l-1},b_l)$对应相对功率$\kappa_l(b_0=0,b_L=1,\kappa_1\ge\cdots\ge\kappa_L>0)$,
实际充电功率为$\pi_s\kappa_l$~kW.
充电会话$q$由进入电量$\varepsilon_q^{\mathrm{in}}$充至离开电量$\varepsilon_q^{\mathrm{out}}$,
补电量$E_q^{\mathrm{ch}}=\varepsilon_q^{\mathrm{out}}-\varepsilon_q^{\mathrm{in}}$,
充电时长为
\begin{equation}
\Delta_q=\frac{3600 B}{\pi_{s(q)}}\sum_{l=1}^{L}
\frac{(\min\{\varepsilon_q^{\mathrm{out}}/B,b_l\}
-\max\{\varepsilon_q^{\mathrm{in}}/B,b_{l-1}\})_+}{\kappa_l}.
\label{eq:charge_duration}
\end{equation}
会话须整段落在其可行窗口$[\underline a_q,\overline a_q]$内:
\begin{equation}
\underline a_q\le a_q^{\mathrm{ch}},\quad a_q^{\mathrm{ch}}+\Delta_q\le\overline a_q.
\label{eq:charge_window}
\end{equation}
会话与电网时段$t$的重叠时长为
\begin{equation}
g_{qt}=[\min\{a_q^{\mathrm{ch}}+\Delta_q,\overline T_t\}-\max\{a_q^{\mathrm{ch}},\underline T_t\}]_+ .
\label{eq:charge_overlap}
\end{equation}
各时段补能量$y_{qt}$由时段内充电曲线积分确定, 满足$\sum_{t\in T}y_{qt}=E_q^{\mathrm{ch}}$;
恒功率($L=1,\kappa_1=1$)退化为$y_{qt}=\pi_{s(q)}g_{qt}/3600$.

充电间接排放按补电量所在时段的电网碳强度核算.
决策时使用预测碳强度$\widehat\gamma_t$, 执行后按结算碳强度$\gamma_t$复算,
模式$p$的预测碳排放与结算碳排放分别为
\begin{equation}
\begin{aligned}
\widehat E_{kp}&=\lambda^g\sum_{r\in\mathcal R_{kp}}\sum_{(i,j)\in r}f_{ij}
+\sum_{q\in\mathcal Q_{kp}}\sum_{t\in T}\widehat\gamma_t y_{qt},\\
E_{kp}&=\lambda^g\sum_{r\in\mathcal R_{kp}}\sum_{(i,j)\in r}f_{ij}
+\sum_{q\in\mathcal Q_{kp}}\sum_{t\in T}\gamma_t y_{qt}.
\end{aligned}
\label{eq:carbon}
\end{equation}
第一项是燃油车直接排放, 第二项是电动车充电侧间接排放.

\subsection{目标函数}

\subsubsection{碳排放成本\texorpdfstring{$F_1$}{F1}}
碳排放成本由燃油车行驶的直接排放和电动车充电的间接排放构成.
燃油油耗由公式(\ref{eq:power})和(\ref{eq:fuel})求得,
规划阶段的充电间接排放按充电发生时段的预测碳强度逐时段核算,
模式$p$的预测碳排放$\widehat E_{kp}$如公式(\ref{eq:carbon})所示.
在碳交易机制中, 政府指定碳配额,
企业碳排放量超过配额$CE$时需要购买更多的碳排放权,
相反也可以通过出售碳排放权赚取利润, 此时碳排放成本的计算如公式(\ref{eq:F1})所示.
\begin{equation}
F_1=p^{\mathrm{car}}\Bigl(\sum_{k\in K}\sum_{p\in\Omega_k}\widehat E_{kp}z_{kp}-CE\Bigr).
\label{eq:F1}
\end{equation}

\subsubsection{固定成本\texorpdfstring{$F_2$}{F2}}
固定成本主要包括车辆折旧费、维修费、保险费和驾驶员工资等, 与派遣的配送趟数成正比.
\begin{equation}
F_2=c^{\mathrm{fix}}\sum_{k\in K}\sum_{p\in\Omega_k}n_{kp}z_{kp}.
\label{eq:F2}
\end{equation}

\subsubsection{里程成本\texorpdfstring{$F_3$}{F3}}
里程成本与行驶距离成正比, 反映轮胎磨损、保养等非能源行驶费用.
\begin{equation}
F_3=\sum_{k\in K}\sum_{p\in\Omega_k}c_k^{\mathrm{km}}D_{kp}z_{kp}.
\label{eq:F3}
\end{equation}

\subsubsection{油耗成本\texorpdfstring{$F_4$}{F4}}
油耗成本与燃油车运输时产生的油耗成正比, 油耗量由能耗模型逐弧累计.
\begin{equation}
F_4=p^f\sum_{k\in K}\sum_{p\in\Omega_k}F_{kp}z_{kp}.
\label{eq:F4}
\end{equation}

\subsubsection{充电成本\texorpdfstring{$F_5$}{F5}}
充电成本包括车场充电电费、公共站充电电费和公共站占用成本.
车场执行分时目录电价, 公共站电价为含服务费的分时总价, 公共站另计占用时间费.
充电发生的时段不同, 电费也随之不同, 因此充电成本按会话在各时段的补电量逐时段结算.
\begin{equation}
F_5=\sum_{k\in K}\sum_{p\in\Omega_k}\Bigl(\sum_{q\in\mathcal Q_{kp}}\sum_{t\in T}p_{s(q),t}^{e}y_{qt}
+c^{\mathrm{occ}}H_{kp}^{S}\Bigr)z_{kp}.
\label{eq:F5}
\end{equation}
其中$s(q)$为会话$q$所在的车场或充电站.

\subsubsection{跨场成本\texorpdfstring{$F_6$}{F6}}
车辆服务非所属车场的客户时计入跨场协调成本,
反映跨场协同产生的信息交换和管理开销.
\begin{equation}
F_6=c^{\mathrm{tr}}\sum_{k\in K}\sum_{p\in\Omega_k}\sum_{i\in N}A_{ikp}\mathbf 1_{d(k)\neq d_i^0}\,z_{kp}.
\label{eq:F6}
\end{equation}

上述公式中, $n_{kp}$, $D_{kp}$, $F_{kp}$, $H_{kp}^{S}$, $\widehat E_{kp}$
分别为模式$p$的配送趟数、总里程(km)、总油耗(L)、公共站占用时间(min)和预测碳排放(kgCO$_2$e),
由模式$p$所含配送趟和充电会话按公式(\ref{eq:power})--(\ref{eq:carbon})逐项累计得到.
$E_{kp}$为同一模式执行后按结算碳强度复算的碳排放.

\subsection{模型建立}

模型分两层建立. 第一层为配送趟约束, 对趟$r$, $d_r^+,d_r^-$为起讫车场副本; 第二层为模式选择约束, $M$为足够大的正数.
\begin{equation}
\min F=F_1+F_2+F_3+F_4+F_5+F_6. \label{eq:obj}
\end{equation}
\begin{equation}
\sum_{j\in V_r}x_{d_r^+jr}=1,\ \sum_{j\in V_r}x_{jd_r^-r}=1,\ \sum_{i\in V_r}x_{id_r^+r}=0,\ \sum_{i\in V_r}x_{d_r^-ir}=0.
\label{eq:trip_flow}
\end{equation}
\begin{equation}
\sum_{j\in V_r}x_{ijr}=\sum_{j\in V_r}x_{jir}=a_{ir},\quad i\in V_r^c.
\label{eq:flow_cons}
\end{equation}
\begin{equation}
L_{d_r^+r}=\sum_{i\in V_r^c}q_i a_{ir},\quad 0\le L_{ir}\le Q_k,\quad i\in V_r.
\label{eq:load_bounds}
\end{equation}
\begin{equation}
L_{jr}\le L_{ir}-q_j+M(1-x_{ijr}),\quad i,j\in V_r.
\label{eq:load_prop}
\end{equation}
\begin{equation}
\underline t_i a_{ir}\le T_{ir}\le\overline t_i+M(1-a_{ir}),\quad i\in V_r^c.
\label{eq:tw}
\end{equation}
\begin{equation}
T_{jr}\ge T_{ir}+\sigma_i+t_{ij}-M(1-x_{ijr}),\quad i,j\in V_r.
\label{eq:time_prop}
\end{equation}
\begin{equation}
B^{\min}\le\varepsilon_{ir}\le B,\quad
\varepsilon_{jr}\le\varepsilon_{ir}+Y_{ir}-e_{ij}+M(1-x_{ijr}),\quad i,j\in V_r.
\label{eq:batt_prop}
\end{equation}
\begin{equation}
Y_{ir}=0,\quad i\in V_r^c.
\label{eq:no_customer_charge}
\end{equation}
\begin{equation}
T_{d_{r'}^+r'}\ge T_{d_r^-r},\quad
\varepsilon_{d_{r'}^+r'}=\varepsilon_{d_r^-r}+\sum_{q\in\mathcal Q_{kp}^{rr'}}E_q^{\mathrm{ch}},\quad r\to r'\ \text{相邻}.
\label{eq:trip_linkage}
\end{equation}
\begin{equation}
\sum_{k\in K}\sum_{p\in\Omega_k}A_{ikp}z_{kp}=1,\quad i\in N,\quad
\sum_{p\in\Omega_k}z_{kp}\le1,\quad k\in K.
\label{eq:coverage}
\end{equation}
\begin{equation}
\sum_{k\in K}\sum_{p\in\Omega_k}O_{skpt}z_{kp}\le C_s,\quad s\in S\cup D,\;t\in T.
\label{eq:station_cap}
\end{equation}
规划阶段的车场收益采用预测碳排放$\widehat E_{kp}$; 执行后的排放报告采用$E_{kp}$.
\begin{equation}
\Pi_d^{\mathrm{stage}}=\sum_{k\in K_d}\sum_{p\in\Omega_k}
\Bigl(\rho\sum_{i\in N}q_iA_{ikp}
-c^{\mathrm{fix}}n_{kp}-c_k^{\mathrm{km}}D_{kp}-p^fF_{kp}-G_{kp}^{e}
-c^{\mathrm{tr}}n_{kp}^{\mathrm{tr}}-p^{\mathrm{car}}\widehat E_{kp}\Bigr)z_{kp}.
\label{eq:stage_profit}
\end{equation}
\begin{equation}
z_{kp}\in\{0,1\},\quad k\in K,\;p\in\Omega_k.
\label{eq:binary}
\end{equation}

公式(\ref{eq:obj})表示运营成本与碳结算成本之和最小化.
公式(\ref{eq:trip_flow})和(\ref{eq:flow_cons})定义趟的起讫车场与客户流守恒;
公式(\ref{eq:load_bounds})和(\ref{eq:load_prop})保证载重不超容量且随服务递减;
公式(\ref{eq:tw})和(\ref{eq:time_prop})保证服务时刻满足硬时间窗并沿弧递推;
公式(\ref{eq:batt_prop})保证电量在允许范围且逐弧传播, 公式(\ref{eq:no_customer_charge})规定客户点不可充电.
公式(\ref{eq:trip_linkage})保证同一实体车相邻趟的时间衔接与跨趟电量连续, 即趟间按需补电,
其中$\mathcal Q_{kp}^{rr'}$为安排在两趟之间的车场充电会话集合;
充电会话窗口由车辆实际停留区间界定, 站内为到离时刻, 趟间为$[T_{d_r^-r},T_{d_{r'}^+r'}]$,
首趟前为运营开始时刻到首趟发车时刻, 所有会话服从公式(\ref{eq:charge_duration})--(\ref{eq:charge_overlap})且互不重叠.
记$\Omega_k$为满足公式(\ref{eq:trip_flow})--(\ref{eq:trip_linkage})的车辆$k$全部可行排班模式的集合,
模式$p\in\Omega_k$完整给定其配送趟、节点时刻、车型和充电会话, 第二层在$\Omega_k$上进行模式选择.
公式(\ref{eq:coverage})保证每个客户由恰好一个模式覆盖且每车至多执行一个模式;
公式(\ref{eq:station_cap})参照容量受限充电站研究\cite{ref:12,ref:13},
保证任意时段站点并发充电数不超过桩数, $O_{skpt}$由会话开始时刻和时长预先计算.
公式(\ref{eq:stage_profit})按实际执行车辆的所属车场计算转移前运营收益,
其中$G_{kp}^{e}$为模式$p$的充电总成本(即公式(\ref{eq:F5})的模式项),
$n_{kp}^{\mathrm{tr}}=\sum_{i\in N}A_{ikp}\mathbf 1_{d(k)\neq d_i^0}$为跨场服务客户数.
基准收益$\Pi_d^0$由同一算法在同等评价预算下,
按车场$d$仅服务自身责任客户、禁用跨场服务的独立经营情形求得.
$\bar\Pi_d^\tau$为截至上次重规划的累计收益, 静态情形取0, 其滚动累计规则见2.6节.

令$D$为参与协作的配送企业集合, 每家企业拥有一个车场及其车队, 并对一部分客户负有配送责任.
对任意非空联盟$A\subseteq D$, 限定只有联盟成员的车场、车辆和责任客户参与配送,
由同一算法求得联盟成本$C(A)$与收入$R(A)$, 并定义可转移收益
\begin{equation}
v(A)=R(A)-C(A),\qquad v(\varnothing)=0.
\label{eq:coalition_profit}
\end{equation}
$|D|$家企业共有$2^{|D|}-1$个非空联盟. 车场$d$的Shapley收益分配为
\begin{equation}
\phi_d=\sum_{A\subseteq D\setminus\{d\}}
\frac{|A|!(|D|-|A|-1)!}{|D|!}
\bigl[v(A\cup\{d\})-v(A)\bigr].
\label{eq:shapley}
\end{equation}
若分配$x$满足
\begin{equation}
\sum_{d\in D}x_d=v(D),\qquad
\sum_{d\in A}x_d\ge v(A),\quad \varnothing\ne A\subset D,
\label{eq:core}
\end{equation}
则没有任何单家或子联盟能通过退出四家合作获得更高收益, 该分配位于核内.
记联盟不满意额为$e(A,x)=v(A)-\sum_{d\in A}x_d$;
核仁按降序对全部$e(A,x)$作字典序最小化, 得到唯一分配$x^{\nu}$.
本文采用核仁作为结算结果, 同时报出Shapley值以比较贡献分配与联盟稳定性的差异.
第二层为集合划分结构, 模式集$\Omega_k$随客户规模指数增长,
本模型是标准VRP的推广, 属于NP-难问题, 第3节给出其启发式求解算法.

\subsection{动态需求处理}

参照动态需求批处理研究\cite{ref:18,ref:71},
本文同时考虑定时和定量两种批事件触发策略.
设第$m-1$次重规划发生于时刻$\tau_{m-1}$,
其后新到事件的累计需求变动量为$W(\tau_{m-1},u)$,
则第$m$次重规划时刻为
\begin{equation}
\tau_m=\min\bigl\{\min\{u:W(\tau_{m-1},u)\ge\overline W\},\;\tau_{m-1}+T^{\mathrm{int}}\bigr\},
\label{eq:trigger}
\end{equation}
即累计需求变动达到上限$\overline W$或距上次重规划达到间隔$T^{\mathrm{int}}$时, 即触发滚动重规划.
在重规划时刻$\tau$, 有效客户集合按新增、取消和已完成服务更新:
\begin{equation}
N^\tau=(N^{\tau-1}\setminus N_{\mathrm{ser}}^\tau\setminus N_{\mathrm{can}}^\tau)\cup N_{\mathrm{new}}^\tau.
\label{eq:customer_update}
\end{equation}
每辆实体车$k$继承当前位置$o_k^\tau$、最早可用时刻$t_k^\tau$、
剩余载重$\ell_k^\tau$和剩余电量$\varepsilon_k^\tau$,
已经完成、已经发车和正在执行的路径前缀不可撤回,
已完成部分的成本、排放和收益按阶段累计:
\begin{equation}
\bar Z^\tau=\bar Z^{\tau-1}+Z^{\tau-1},\quad
\bar E^\tau=\bar E^{\tau-1}+\Delta E^{\tau-1},\quad
\bar\Pi_d^\tau=\bar\Pi_d^{\tau-1}+\Pi_d^{\mathrm{stage},\tau-1}.
\label{eq:accumulation}
\end{equation}
状态继承保证滚动重规划不撤回已执行路径、不重复使用实体车、不重复结算历史排放.

% ================================================================
% SECTION 3: 算法设计 (按陈雨蝶: 流程图总述段+叙述式步骤+示意图)
% ================================================================

\section{算法设计}

\subsection{算法流程图}

%% MATH_REF: Vidal 2012/2014 (HGS); Rochat & Taillard 1995 (路线池); Wouda 2024 (PyVRP)
上述模型同时包含跨场服务、车型选择、实体车多趟、分时充电排放、成员结算参与和动态订单,
其可行解需同时满足路径层与实体车排班层的耦合约束.
在该结构下, 仅凭路线邻域无法判定一个候选方案是否可行:
可行与否取决于该方案被补全为实体车排班并完成补能之后的结果.
因此本文不另行设计搜索内核, 而是直接采用成熟的混合遗传搜索
(Hybrid Genetic Search, HGS)\cite{ref:76,ref:77}作为搜索底座,
并在其上构造面向本问题的特化求解组件:
统一完整评价、实体车多趟补全、碳感知充电排程与动态状态继承.
算法记为HGS-FRVCP（融合固定路线充电求解的混合遗传搜索）, 主流程如图\ref{fig:algorithm-flow}所示.
算法按代循环推进: 每代由交叉产生子代, 子代依次经过串行改进链的三步问题特化改进,
再由完整评价器裁决能否入群, 如图\ref{fig:route-pool}所示.
图\ref{fig:algorithm-flow}给出总流程, 与下文步骤1--6逐项对应.
算法全过程坚持同一核算口径: 候选方案能否被接受、能否进入种群、能否存活,
均由同一完整评价器给出的完整成本与完整可行性决定,
不以简化路线代价近似替代完整方案代价.
能源消耗、补能安排、趟间时间与电量衔接、分时电价和时变碳强度进入同一本账;
候选补全失败时保留失败原因, 不以初始解顶替.
第4节报告其求解结果。

\begin{figure}[H]
  \centering
  \includegraphics{generated_figures/figure_1_algorithm_flow.pdf}
  \caption{本文算法流程图}
  \label{fig:algorithm-flow}
\end{figure}

\subsection{算法步骤}

\textbf{步骤1：读取状态并构造可行起点.} 静态阶段读取订单、车场、有限实体车队与分时电价、逐时电网碳强度序列；动态阶段读取式\eqref{eq:customer_update}--\eqref{eq:accumulation}定义的剩余客户与车辆状态。初始个体先由完整模型检查，缺客户或未通过完整评价的候选不进入种群。

\textbf{步骤2：执行混合遗传搜索主循环.} 算法维护可行与不可行两个子种群，按锦标赛选择父代并执行交叉。罚系数依可行解比例自适应调整。停止条件与评价预算由本项目在同一算例、同一机器上的实测收敛记录标定，不直接照搬其他算法的迭代数。

\textbf{步骤3：分阶段提名与最优改进局部搜索.} 交叉子代进入局部搜索，其候选动作由提名器按阶段顺序供给，如图\ref{fig:route-pool}所示：先给出路线邻域动作，包括客户重定位、客户交换、路段反转、路段重定位与尾段交换；再给出机制动作，包括整车车型交换、充电时刻调整与整趟交换。两个阶段的动作按顺序拼接并去重，汇入同一条提名流。局部搜索每轮扫描该提名流的全部候选，按完整评价结果选出其中最优的一个动作并应用，重复至本轮不再产生改善。提名阶段的先后只决定候选出现的次序，取舍一律由完整评价器裁决，不以代理分数替代完整方案代价；每代只将一个最终子代送入种群。

\begin{figure}[H]
  \centering
  \includegraphics{generated_figures/figure_2_serial_chain.pdf}
  \caption{分阶段提名与最优改进局部搜索}
  \label{fig:route-pool}
\end{figure}

\textbf{步骤4：完成实体车排班与补能.} 将配送趟指派给具体实体车，按时间顺序衔接多趟任务，并在首趟出发前与相邻趟之间补齐电量。补能过程按分段充电功率函数计算时长，充电电费与充电排放分别按该时段的分时电价与电网碳强度结算。违反时间窗、载重、电量、公共站容量或实体车上限的候选保留失败原因，不以初始解顶替。

\textbf{步骤5：完整评价、入群与生存选择.} 候选方案的完整成本与完整可行性由同一评价器给出；入群按完整可行性分池，生存选择与父代选择使用同一本账。全局最好解必须完整可行，且仅在严格改善时更新。

\textbf{步骤6：独立验解与动态调用.} 进入种群的候选重新执行一次独立检查，避免评价与搜索共用同一实现带来的系统性偏差。动态事件触发后，冻结已经执行的路径前缀，继承各实体车的位置、最早可用时刻、剩余载重与剩余电量，并把已完成部分的成本与排放按阶段累计，随后仅对剩余任务返回步骤2重新求解。


\section{数值实验}

\subsection{算例与参数}

代表实例为京津冀50客户、2车场、97个公共站节点的构造算例，客户责任按编号等分为两家企业各25个。动态事件流包含10个新增客户事件，接单窗为08:00--10:00，共5个触发批；分时电价和碳强度的信息分辨率为1小时。

车辆成本、容量与柴油价格取自本文算例的车辆参数与当地价格表；电动车非能源里程成本中的电池折旧部分按电池价格与寿命里程折算，寿命里程参照Goeke与Schneider\cite{ref:27}，电动车相对燃油车每辆每日50元的固定成本溢价参照陈婉茹等\cite{ref:23}；车场单车额定充电功率参照中国地方新能源汽车充（换）电设施专项规划；分时电价与逐时电网碳强度取自北京2025年2月12日的当地价格表与省级逐小时碳强度投影数据集，信息分辨率为1小时。尚未闭合到原始页码的出处为\TBD{出处或页码}。关键参数如表\ref{tab:parameters}所示。

\begin{table}[htbp]
  \centering
  \caption{关键车辆、充电、电价与碳强度参数及来源}
  \label{tab:parameters}
  \scriptsize
  \begin{tabular}{lrl}
    \toprule
    参数 & 数值 & 单位\\
    \midrule
    燃油车日固定成本 & 170.00 & 元/(辆·日)\\
    电动车日固定成本 & 220.00 & 元/(辆·日)\\
    燃油车非能源里程成本 & 0.7800 & 元/km\\
    电动车非能源里程成本 & 0.9145 & 元/km\\
    车辆容积容量 & 7.2 & m$^3$\\
    柴油价格 & 7.48 & 元/L\\
    车场单车额定充电功率 & 60.0 & kW\\
    车场分时电价（谷/平/峰） & 0.5633／0.8364／1.1486 & 元/kWh\\
    公共站服务费 & 0.4 & 元/kWh\\
    公共站分时总电价（谷/平/峰） & 0.9633／1.2364／1.5486 & 元/kWh\\
    逐时电网碳强度（24 个小时值） & 0.1541～0.6439 & kgCO$_2$e/kWh\\
    单位碳价 & 0.07502 & 元/kgCO$_2$e\\
    \bottomrule
  \end{tabular}
\end{table}

\subsection{公开算例上的求解结果}

公开实验采用28个多车场带时间窗算例检验基础路径搜索能力。当前每题只有一次运行，公开最好值及其版本、已发表算法逐题值和多种子均值尚未补齐。28题均完成全部客户并通过原始精度路线复核；完成需求量/总需求量仍为\TBD{逐题服务需求字段}。因此本节只登记逐题单次值，不使用“优于”“领先”或“胜出”。

\begin{table}[htbp]
  \centering
  \caption{公开算例上不同算法的求解结果}
  \label{tab:public28}
  \scriptsize
  \resizebox{\textwidth}{!}{%
  \begin{tabular}{lrrrrrrrrr}
    \toprule
    & & \multicolumn{2}{c}{VCGP\cite{ref:74}} & \multicolumn{2}{c}{MDFIHA\cite{ref:75}} & \multicolumn{2}{c}{MDFIHA-ETGA\cite{ref:75}} & \multicolumn{2}{c}{本文算法}\\
    \cmidrule(lr){3-4}\cmidrule(lr){5-6}\cmidrule(lr){7-8}\cmidrule(lr){9-10}
    算例 & BKS & Best & Avg & Best & Avg & Best & Avg & Best & Avg\\
    \midrule
    PR11A & \TBD{BKS} & \TBD{逐题值} & \TBD{均值} & \TBD{逐题值} & \TBD{均值} & \TBD{逐题值} & \TBD{均值} & 6776.761 & \TBD{多种子}\\
    PR11B & \TBD{BKS} & \TBD{逐题值} & \TBD{均值} & \TBD{逐题值} & \TBD{均值} & \TBD{逐题值} & \TBD{均值} & 4820.376 & \TBD{多种子}\\
    PR12A & \TBD{BKS} & \TBD{逐题值} & \TBD{均值} & \TBD{逐题值} & \TBD{均值} & \TBD{逐题值} & \TBD{均值} & 8251.671 & \TBD{多种子}\\
    PR12B & \TBD{BKS} & \TBD{逐题值} & \TBD{均值} & \TBD{逐题值} & \TBD{均值} & \TBD{逐题值} & \TBD{均值} & 6041.261 & \TBD{多种子}\\
    PR13A & \TBD{BKS} & \TBD{逐题值} & \TBD{均值} & \TBD{逐题值} & \TBD{均值} & \TBD{逐题值} & \TBD{均值} & 9722.412 & \TBD{多种子}\\
    PR13B & \TBD{BKS} & \TBD{逐题值} & \TBD{均值} & \TBD{逐题值} & \TBD{均值} & \TBD{逐题值} & \TBD{均值} & 7293.739 & \TBD{多种子}\\
    PR14A & \TBD{BKS} & \TBD{逐题值} & \TBD{均值} & \TBD{逐题值} & \TBD{均值} & \TBD{逐题值} & \TBD{均值} & 11280.749 & \TBD{多种子}\\
    PR14B & \TBD{BKS} & \TBD{逐题值} & \TBD{均值} & \TBD{逐题值} & \TBD{均值} & \TBD{逐题值} & \TBD{均值} & 8825.714 & \TBD{多种子}\\
    PR15A & \TBD{BKS} & \TBD{逐题值} & \TBD{均值} & \TBD{逐题值} & \TBD{均值} & \TBD{逐题值} & \TBD{均值} & 13229.546 & \TBD{多种子}\\
    PR15B & \TBD{BKS} & \TBD{逐题值} & \TBD{均值} & \TBD{逐题值} & \TBD{均值} & \TBD{逐题值} & \TBD{均值} & 10522.498 & \TBD{多种子}\\
    PR16A & \TBD{BKS} & \TBD{逐题值} & \TBD{均值} & \TBD{逐题值} & \TBD{均值} & \TBD{逐题值} & \TBD{均值} & 14633.726 & \TBD{多种子}\\
    PR16B & \TBD{BKS} & \TBD{逐题值} & \TBD{均值} & \TBD{逐题值} & \TBD{均值} & \TBD{逐题值} & \TBD{均值} & 11762.652 & \TBD{多种子}\\
    PR17A & \TBD{BKS} & \TBD{逐题值} & \TBD{均值} & \TBD{逐题值} & \TBD{均值} & \TBD{逐题值} & \TBD{均值} & 6316.228 & \TBD{多种子}\\
    PR17B & \TBD{BKS} & \TBD{逐题值} & \TBD{均值} & \TBD{逐题值} & \TBD{均值} & \TBD{逐题值} & \TBD{均值} & 4796.924 & \TBD{多种子}\\
    PR18A & \TBD{BKS} & \TBD{逐题值} & \TBD{均值} & \TBD{逐题值} & \TBD{均值} & \TBD{逐题值} & \TBD{均值} & 8365.857 & \TBD{多种子}\\
    PR18B & \TBD{BKS} & \TBD{逐题值} & \TBD{均值} & \TBD{逐题值} & \TBD{均值} & \TBD{逐题值} & \TBD{均值} & 6512.75 & \TBD{多种子}\\
    PR19A & \TBD{BKS} & \TBD{逐题值} & \TBD{均值} & \TBD{逐题值} & \TBD{均值} & \TBD{逐题值} & \TBD{均值} & 10831.511 & \TBD{多种子}\\
    PR19B & \TBD{BKS} & \TBD{逐题值} & \TBD{均值} & \TBD{逐题值} & \TBD{均值} & \TBD{逐题值} & \TBD{均值} & 8199.189 & \TBD{多种子}\\
    PR20A & \TBD{BKS} & \TBD{逐题值} & \TBD{均值} & \TBD{逐题值} & \TBD{均值} & \TBD{逐题值} & \TBD{均值} & 12204.194 & \TBD{多种子}\\
    PR20B & \TBD{BKS} & \TBD{逐题值} & \TBD{均值} & \TBD{逐题值} & \TBD{均值} & \TBD{逐题值} & \TBD{均值} & 10324.916 & \TBD{多种子}\\
    PR21A & \TBD{BKS} & \TBD{逐题值} & \TBD{均值} & \TBD{逐题值} & \TBD{均值} & \TBD{逐题值} & \TBD{均值} & 6251.424 & \TBD{多种子}\\
    PR21B & \TBD{BKS} & \TBD{逐题值} & \TBD{均值} & \TBD{逐题值} & \TBD{均值} & \TBD{逐题值} & \TBD{均值} & 4893.444 & \TBD{多种子}\\
    PR22A & \TBD{BKS} & \TBD{逐题值} & \TBD{均值} & \TBD{逐题值} & \TBD{均值} & \TBD{逐题值} & \TBD{均值} & 8121.877 & \TBD{多种子}\\
    PR22B & \TBD{BKS} & \TBD{逐题值} & \TBD{均值} & \TBD{逐题值} & \TBD{均值} & \TBD{逐题值} & \TBD{均值} & 6516.749 & \TBD{多种子}\\
    PR23A & \TBD{BKS} & \TBD{逐题值} & \TBD{均值} & \TBD{逐题值} & \TBD{均值} & \TBD{逐题值} & \TBD{均值} & 10138.318 & \TBD{多种子}\\
    PR23B & \TBD{BKS} & \TBD{逐题值} & \TBD{均值} & \TBD{逐题值} & \TBD{均值} & \TBD{逐题值} & \TBD{均值} & 8480.755 & \TBD{多种子}\\
    PR24A & \TBD{BKS} & \TBD{逐题值} & \TBD{均值} & \TBD{逐题值} & \TBD{均值} & \TBD{逐题值} & \TBD{均值} & 12149.427 & \TBD{多种子}\\
    PR24B & \TBD{BKS} & \TBD{逐题值} & \TBD{均值} & \TBD{逐题值} & \TBD{均值} & \TBD{逐题值} & \TBD{均值} & 10905.516 & \TBD{多种子}\\
    \midrule
    Average & \TBD{平均BKS} & \TBD{平均值} & \TBD{平均值} & \TBD{平均值} & \TBD{平均值} & \TBD{平均值} & \TBD{平均值} & \TBD{平均值} & \TBD{多种子}\\
    Nbks & 28 & \TBD{题数} & \TBD{题数} & \TBD{题数} & \TBD{题数} & \TBD{题数} & \TBD{题数} & \TBD{题数} & \TBD{题数}\\
    \bottomrule
  \end{tabular}
  }
\end{table}

\subsection{代表实例最终解}

表\ref{tab:final-solution}把一辆实体车的一日多趟拆成多行，车辆号和车辆固定成本只在该车第一趟出现。逐趟客户序列已由保存解登记；发车/返回时刻没有原始字段，继续保留占位；逐车日里程只在首行列一次。

\begin{table}[htbp]
  \centering
  \caption{代表实例最终配送方案}
  \label{tab:final-solution}
  \scriptsize
  \resizebox{\textwidth}{!}{%
  \begin{tabular}{llllllrr}
    \toprule
    实体车 & 所属车场 & 车型 & 趟次 & 客户序列 & 发车/返回时刻 & 本车日里程（km） & 车辆固定成本（元）\\
    \midrule
    CV\_A\_1 & 车场A & 燃油车 & 1 & C035--C022 & \TBD{逐趟时刻} & 151.083 & 170.000\\
              &       &        & 2 & C014--C028--C006--C026 & \TBD{逐趟时刻} & & \\
    CV\_A\_2 & 车场A & 燃油车 & 1 & C031--C017 & \TBD{逐趟时刻} & 83.613 & 170.000\\
    CV\_B\_1 & 车场B & 燃油车 & 1 & C024--C023--C019--C029 & \TBD{逐趟时刻} & 56.999 & 170.000\\
    CV\_B\_2 & 车场B & 燃油车 & 1 & C025--C015--C040--C013 & \TBD{逐趟时刻} & 151.064 & 170.000\\
              &       &        & 2 & C034--C047--C049 & \TBD{逐趟时刻} & & \\
              &       &        & 3 & C048--C020--C037 & \TBD{逐趟时刻} & & \\
    CV\_B\_3 & 车场B & 燃油车 & 1 & C033--C043--C038 & \TBD{逐趟时刻} & 162.811 & 170.000\\
              &       &        & 2 & C036--C030--C041 & \TBD{逐趟时刻} & & \\
              &       &        & 3 & C009--C007 & \TBD{逐趟时刻} & & \\
    CV\_B\_4 & 车场B & 燃油车 & 1 & C002--C012 & \TBD{逐趟时刻} & 93.774 & 170.000\\
              &       &        & 2 & C008--C050--C046--C021 & \TBD{逐趟时刻} & & \\
    CV\_B\_5 & 车场B & 燃油车 & 1 & C001--C027--C016 & \TBD{逐趟时刻} & 119.043 & 170.000\\
              &       &        & 2 & C005--C045--C018--C039 & \TBD{逐趟时刻} & & \\
    CV\_B\_6 & 车场B & 燃油车 & 1 & C010--C044--C011 & \TBD{逐趟时刻} & 117.142 & 170.000\\
              &       &        & 2 & C042--C004--C032--C003 & \TBD{逐趟时刻} & & \\
    \bottomrule
  \end{tabular}
  }
\end{table}

（1）该保存解完成50/50个客户和13264/13264~kg需求，使用8辆实体车执行16趟，总里程935.530~km。车辆固定成本1360.000元按8辆实体车各计一次，不按16趟重复计取。

（2）单目标总账为3033.206元，其中固定成本1360.000元、里程成本729.713元、燃油成本919.139元、电费0.000元、碳价折算成本24.354元。

（3）燃油消耗为122.879~L，燃油车直接排放324.636~kgCO$_2$e；该解没有派遣电动车，因此电能消耗和电动车充电排放均为0。

\subsection{本文算法收敛过程}

图\ref{fig:convergence}的曲线来自统一算例联合臂的一次探针运行，共12个过程点，横轴为实际运行时间。该运行未通过整包验收，图中数据按占位登记，不作为正式收敛结论；正式收敛标定须在算法定稿后按本项目实测轨迹重做。

\begin{figure}[htbp]
  \centering
  \includegraphics{generated_figures/figure_2_algorithm_convergence.pdf}
  \caption{不同算法迭代图}
  \label{fig:convergence}
\end{figure}

\subsection{混合车队车型指派层}

受购置成本与补能条件限制, 车队无法一次性完成电动化替换,
燃油车与电动车在相当长时期内共存.
企业由此面对的第一个决策问题是: 两类车型的续航、能耗与成本结构不同,
何种配送任务应交由电动车、何种应保留给燃油车,
以及车型构成改变之后, 原有的最优路径是否仍然最优.
本节在同一实例上给定不同的车型可用名额, 比较实际派遣构成、路径结构与成本排放.

车型层登记18CV/2EV、13CV/7EV、7CV/13EV和2CV/18EV四个可用名额档。前两档完成探针，均完成50/50个客户和13264/13264~kg需求，但实际派遣电动车为0；后两档因完成解引用未登记实体车而失败，不能写成模型不可行。

\begin{table}[htbp]
  \centering
  \caption{混合车队可用名额与实际派遣对比}
  \label{tab:fleet-levels}
  \scriptsize
  \resizebox{\textwidth}{!}{%
  \begin{tabular}{lrrrr}
    \toprule
    指标（单位） & 18CV/2EV & 13CV/7EV & 7CV/13EV & 2CV/18EV\\
    \midrule
    单目标总账（元） & \textbf{3339.410} & 3372.864 & \TBD{搜索未完成} & \TBD{搜索未完成}\\
    车辆固定成本（元） & \textbf{1530.000} & 1700.000 & \TBD{搜索未完成} & \TBD{搜索未完成}\\
    非能源里程成本（元） & 790.240 & \textbf{730.169} & \TBD{搜索未完成} & \TBD{搜索未完成}\\
    燃油成本（元） & 992.863 & \textbf{918.361} & \TBD{搜索未完成} & \TBD{搜索未完成}\\
    电费（元） & \textbf{0.000} & \textbf{0.000} & \TBD{搜索未完成} & \TBD{搜索未完成}\\
    碳价折算成本（元） & 26.308 & \textbf{24.334} & \TBD{搜索未完成} & \TBD{搜索未完成}\\
    可用燃油车/电动车（辆） & 18/2 & 13/7 & 7/13 & 2/18\\
    实际派遣燃油车/电动车（辆） & 9/0 & 10/0 & \TBD{搜索未完成} & \TBD{搜索未完成}\\
    配送趟次（趟） & 17 & \textbf{16} & \TBD{搜索未完成} & \TBD{搜索未完成}\\
    总里程（km） & 1013.128 & \textbf{936.115} & \TBD{搜索未完成} & \TBD{搜索未完成}\\
    燃油车直接排放（kgCO$_2$e） & 350.675 & \textbf{324.361} & \TBD{搜索未完成} & \TBD{搜索未完成}\\
    电动车充电排放（kgCO$_2$e） & \textbf{0.000} & \textbf{0.000} & \TBD{搜索未完成} & \TBD{搜索未完成}\\
    \bottomrule
  \end{tabular}
  }
\end{table}

\begin{figure}[htbp]
  \centering
  \includegraphics{generated_figures/figure_4_mixed_fleet.pdf}
  \caption{不同碳价下混合车队的车型构成}
  \label{fig:fleet-carbon-price}
\end{figure}

图\ref{fig:fleet-carbon-price}的横轴保持约0--2.5元/kg，纵轴为实际车辆数，燃油车与电动车分别成线。本算例已测得的0.449元/kg只作为“见图3”的淡参考线；全电动阈值保持\TBD{本算例实测}。

\subsection{非线性充电对照}

电动车一旦真正投入运行, 补能过程立即成为路径可行性的组成部分.
真实补能并非恒功率: 荷电状态越高, 单位时间充入的电量越少.
若按恒功率刻画补能, 补能时长会被系统性低估,
一条路线能否在时间窗内完成、下一趟能否按时发车都可能被判断错误.
本节以线性恒功率为影子对照, 检验真实充电曲线是否改变车型配置与路径安排.

非线性充电按S/R分组：S表示简化充电函数下的账面求解，R表示同一方案按参考非线性函数回算。当前没有完整S/R正式对照，完成客户数和完成需求量也没有合格数字，表\ref{tab:nonlinear-charging}全部留白。

\begin{table}[htbp]
  \centering
  \caption{不同充电功率函数的账面求解与非线性回算对比}
  \label{tab:nonlinear-charging}
  \scriptsize
  \resizebox{\textwidth}{!}{%
  \begin{tabular}{lrrrrr}
    \toprule
    & \multicolumn{2}{c}{恒功率} & \multicolumn{2}{c}{分段SOC功率} & 非线性SOC功率\\
    \cmidrule(lr){2-3}\cmidrule(lr){4-5}\cmidrule(lr){6-6}
    指标（单位） & S & R & S & R & R\\
    \midrule
    单目标总账（元） & \TBD{正式S/R} & \TBD{正式S/R} & \TBD{正式S/R} & \TBD{正式S/R} & \TBD{正式S/R}\\
    车辆固定成本（元） & \TBD{正式S/R} & \TBD{正式S/R} & \TBD{正式S/R} & \TBD{正式S/R} & \TBD{正式S/R}\\
    非能源里程成本（元） & \TBD{正式S/R} & \TBD{正式S/R} & \TBD{正式S/R} & \TBD{正式S/R} & \TBD{正式S/R}\\
    燃油成本（元） & \TBD{正式S/R} & \TBD{正式S/R} & \TBD{正式S/R} & \TBD{正式S/R} & \TBD{正式S/R}\\
    电费（元） & \TBD{正式S/R} & \TBD{正式S/R} & \TBD{正式S/R} & \TBD{正式S/R} & \TBD{正式S/R}\\
    碳价折算成本（元） & \TBD{正式S/R} & \TBD{正式S/R} & \TBD{正式S/R} & \TBD{正式S/R} & \TBD{正式S/R}\\
    可行运行数（次） & \TBD{正式S/R} & \TBD{正式S/R} & \TBD{正式S/R} & \TBD{正式S/R} & \TBD{正式S/R}\\
    充电会话数（次） & \TBD{正式S/R} & \TBD{正式S/R} & \TBD{正式S/R} & \TBD{正式S/R} & \TBD{正式S/R}\\
    充电量（kWh） & \TBD{正式S/R} & \TBD{正式S/R} & \TBD{正式S/R} & \TBD{正式S/R} & \TBD{正式S/R}\\
    充电时长（min） & \TBD{正式S/R} & \TBD{正式S/R} & \TBD{正式S/R} & \TBD{正式S/R} & \TBD{正式S/R}\\
    会话最高SOC（\%） & \TBD{正式S/R} & \TBD{正式S/R} & \TBD{正式S/R} & \TBD{正式S/R} & \TBD{正式S/R}\\
    充电排放（kgCO$_2$e） & \TBD{正式S/R} & \TBD{正式S/R} & \TBD{正式S/R} & \TBD{正式S/R} & \TBD{正式S/R}\\
    完成客户数（个） & \TBD{正式S/R} & \TBD{正式S/R} & \TBD{正式S/R} & \TBD{正式S/R} & \TBD{正式S/R}\\
    完成需求量（kg） & \TBD{正式S/R} & \TBD{正式S/R} & \TBD{正式S/R} & \TBD{正式S/R} & \TBD{正式S/R}\\
    \bottomrule
  \end{tabular}
  }
\end{table}

\subsection{时变碳强度下的充电时刻层}

补能过程刻画真实之后, 更深一层的问题随之显现:
电动车并非只要用电就天然低碳.
同样充入一度电, 在不同时段对应的电网碳排放并不相同,
因此不能只回答充入多少, 还须回答何时充入.
而何时可充又取决于何时发车、服务哪些客户、何时返场以及在何处补能——
补能窗口本身是配送排班的产物, 而非外生给定.
本节考察车辆是否值得为进入低碳时段而改变自身的运行安排.

固定平均碳强度、时变有空即充、时变电费最省和时变碳感知四个对照臂使用同一实例与服务要求。当前登记值来自固定排班下的精确枚举，电费最省臂尚无结果。已有三臂完成50/50个客户和13264.0/13264.0~kg需求；新增成本分项与完成需求量按探针包原始读数占位展示。五项成本中固定、里程和燃油成本在该固定排班对照中保持不变，合计2877.208798592812元。

\begin{table}[htbp]
  \centering
  \caption{时变电网碳强度与充电选择规则对比}
  \label{tab:carbon-charging}
  \scriptsize
  \resizebox{\textwidth}{!}{%
  \begin{tabular}{lrrrrr}
    \toprule
    指标（单位） & 固定平均碳强度 & 时变有空即充 & 时变电费最省 & 时变碳感知 & 碳感知相对固定基准\\
    \midrule
    单目标总账（元） & \textbf{3032.527} & 3032.970 & \TBD{该策略} & 3055.063 & +0.743\%\\
    车辆固定成本（元） & 1780.000 & 1780.000 & 1780.000 & 1780.000 & 0\\
    非能源里程成本（元） & 811.907 & 811.907 & 811.907 & 811.907 & 0\\
    燃油成本（元） & 285.302 & 285.302 & 285.302 & 285.302 & 0\\
    电费（元） & \textbf{141.525} & \textbf{141.525} & \textbf{\TBD{该策略}} & 168.049 & +18.742\%\\
    碳价折算成本（元，碳价0.075） & 6.236 & 6.679 & \TBD{该策略} & \textbf{2.247} & $-63.961\%$\\
    充电侧总账（元） & \textbf{147.761} & 148.203 & \TBD{该策略} & 170.297 & +15.252\%\\
    充电量（kWh） & 189.279 & 189.279 & \TBD{该策略} & 189.279 & 0\\
    可移动充电量（kWh） & 0.000 & 0.000 & \TBD{该策略} & \textbf{174.732} & ---\\
    充电排放（kgCO$_2$e） & 83.150 & 89.048 & \TBD{该策略} & \textbf{29.967} & $-63.961\%$\\
    配送作业阶段总排放（kgCO$_2$e） & 183.917 & 189.816 & \TBD{该策略} & \textbf{130.734} & $-28.917\%$\\
    完成需求量/总需求量（kg） & 13264.0/13264.0 & 13264.0/13264.0 & \TBD{该策略} & 13264.0/13264.0 & 0\\
    \bottomrule
  \end{tabular}
  }
\end{table}

表\ref{tab:carbon-charging}按充电账、全系统排放账和全系统钱账三个尺度递进呈现。固定平均因子0.4393~kgCO$_2$e/kWh会抹平时刻差异；时变碳感知把174.732~kWh移入更清洁窗口，使充电排放减少59.082~kgCO$_2$e，但在0.075元/kg碳价下电费增加26.525元，完整目标净增22.094元。

\begin{figure}[htbp]
  \centering
  \includegraphics[width=\textwidth]{generated_figures/figure_3_carbon_tariff_charging.pdf}
  \caption{电网碳强度、分时电价与充电时刻}
  \label{fig:carbon-charging}
\end{figure}

图\ref{fig:carbon-charging}(a)显示，电价谷段与电网碳强度低谷并不重合：谷段八小时的碳强度处于全天最高区间，最高值0.644~kgCO$_2$e/kWh出现在凌晨谷段之内，而全天最低值0.154出现在电价平段的午间。两个信号在时间上方向相反，仅按电费最省安排充电，会把充电推向碳强度最高的时段。图\ref{fig:carbon-charging}(b)进一步表明，这一错位在本文的配送排班下确实可被利用：两种规则充入的总电量相同，均为189.3~kWh，但有空即充把其中103.9~kWh置于凌晨谷段，碳感知择时则把143.6~kWh移至午间低碳时段。可移动性来自实体车多趟排班——车辆日间返场装货的间隔恰好落在碳强度低谷区间，因而已登记的10次充电中有9次改期成功，92.3\%的电量得以移动，仅14.5~kWh因窗口不可用而保留原位。该窗口并非外生给定，改变路线与趟次划分即会改变可用的补能时段。

\subsection{碳价情景扫描}

上一节表明, 改变充电时刻可以降低充电环节排放, 但同时抬高电费支出.
在现行碳价水平下, 企业按自身账目并不会主动作出这一改变;
这并非企业不重视减排, 而是该项改变在其成本核算中不成立.
若要使企业的经济决策与减排目标一致, 碳价须提高到何种水平, 是政策制定者必须回答的问题.
本节按阶梯逐点抬高单位碳价, 定位二者对齐的临界区间.

表\ref{tab:carbon-price}独立承载政策情景。每个碳价点对应碳感知规则与有空即充规则的配对完整目标差；正值表示碳感知更贵，负值表示碳感知更省。两策略的完成客户数和完成需求量沿用充电时刻层的同批服务要求，具体需求字段仍为\TBD{配对服务需求量}。

\begin{table}[htbp]
  \centering
  \caption{单位碳价情景扫描结果}
  \label{tab:carbon-price}
  \small
  \begin{tabular}{rr}
    \toprule
    单位碳价（元/kgCO$_2$e） & 碳感知相对有空即充的完整目标变化（元）\\
    \midrule
    0.075 & +22.094\\
    0.20 & +14.708\\
    0.60 & $-8.924$\\
    1.10 & $-38.465$\\
    1.50 & $-62.098$\\
    2.10 & $-97.547$\\
    \bottomrule
  \end{tabular}
\end{table}

配对算式给出的盈亏平衡碳价为0.4489498210731408元/kg。表内0.20元/kg时碳感知仍多花14.708元，0.60元/kg时已少花8.924元；这两个格点把0.449元/kg阈值夹在中间。

\subsection{多车场配送协同}

至此各项优化仍在单个企业的客户责任边界之内进行.
现实中常见的情形是: 一个车场任务饱和而另一个车场相对空闲;
某客户在地理上明显更靠近另一企业的车场, 却因既有归属只能由本场服务.
这说明进一步限制效率的可能已不是车辆与补能能力, 而是人为划定的归属边界.
本节打开该边界, 允许跨企业重新分配客户与任务, 考察由此能否得到更优的配送路径.

协同效率比较两家各自重新优化后的单干成本之和与关闭公平约束后的联合优化成本。两臂均完成50/50个客户和13264/13264~kg需求，使用8辆实体车执行16趟，物理违规为零。已登记结果中，两家单干成本之和为3922.079元，联合成本为3033.206元；表中成本分项及其变化比例按现有探针保存解占位展示。

\begin{table}[htbp]
  \centering
  \caption{多车场配送协同模式对比}
  \label{tab:synergy}
  \scriptsize
  \begin{tabular}{lrrr}
    \toprule
    指标（单位） & 两家单干之和 & 联合协同 & 联合相对单干\\
    \midrule
    单目标总账（元） & 3922.079 & \textbf{3033.206} & $-22.663\%$\\
    车辆固定成本（元） & 1660.000 & 1360.000 & $-18.072\%$\\
    非能源里程成本（元） & 1517.484 & 729.713 & $-51.913\%$\\
    燃油成本（元） & 386.942 & 919.139 & +137.539\%\\
    电费（元） & 333.197 & \textbf{0.000} & $-100.000\%$\\
    碳价折算成本（元） & 24.456 & \textbf{24.354} & $-0.417\%$\\
    启用燃油车/电动车（辆） & 2/6 & 8/0 & 构成变化\\
    启用实体车数（辆） & \textbf{8} & \textbf{8} & 0\\
    配送趟次（趟） & \textbf{16} & \textbf{16} & 0\\
    总里程（km） & 1717.715 & \textbf{935.530} & $-45.536\%$\\
    燃油车直接排放（kgCO$_2$e） & \textbf{136.666} & 324.636 & +137.540\%\\
    电动车充电排放（kgCO$_2$e） & 189.328 & \textbf{0.000} & $-100.000\%$\\
    配送作业阶段总排放（kgCO$_2$e） & 325.994 & \textbf{324.636} & $-0.417\%$\\
    \bottomrule
  \end{tabular}
\end{table}

表\ref{tab:synergy}回答的是整体协同效率，不把成本变化单独归因于跨场路线。协同方案相对两家单干节约888.873元；车型构成同时由2辆燃油车/6辆电动车变为8辆燃油车/0辆电动车，直接排放与充电排放的变化方向相反。

\begin{figure}[htbp]
  \centering
  \includegraphics{generated_figures/concept_collaboration_fairness_notitle.pdf}
  \caption{独立配送与联合配送最终路径图（示意）}
  \label{fig:collaboration-fairness-routes}
\end{figure}

图\ref{fig:collaboration-fairness-routes}以示意方式对照独立配送、联合配送和转归解；具体协同数字见表\ref{tab:synergy}。

\subsection{协同收益公平分配}

协同一旦真正带来收益, 新的矛盾随即出现:
系统整体改善并不意味着每个参与方都改善.
总成本下降的同时, 个别车场可能承担更多任务而自身利益受损;
此类方案即使在数学上最优, 现实中也不会有企业愿意持续合作.
本节在保证各参与方均可接受的前提下考察路径与收益应如何重新安排,
以及为此需要牺牲多少系统效率.

公平分配与协同效率分开：它固定同一联合方案，不重新优化路线，只改变3033.206元联合总成本在两家企业之间的承担额。现有两方数据下，Shapley与核仁给出相同数字；最终采用哪一个方法名仍为\TBD{用户批准的分摊方法}。

\begin{table}[htbp]
  \centering
  \caption{协同收益公平分配结果}
  \label{tab:allocation}
  \small
  \begin{tabular}{llrr}
    \toprule
    分摊方法 & 指标 & 车场A & 车场B\\
    \midrule
    \multirow{4}{*}{\TBD{分摊方法}} & 单干成本（元） & 2068.897 & 1853.181\\
      & 分摊成本（元） & \textbf{1624.461} & \textbf{1408.745}\\
      & 节约量（元） & 444.436 & 444.436\\
      & 节约百分比（\%） & 21.482 & 23.982\\
    \bottomrule
  \end{tabular}
\end{table}

分摊后两家各节约约444.436元，均不高于各自单干成本。路线归属记账和事后分摊是两套成本承担方式，不把协同路线的成本节约与公平分配混写成同一机制。图\ref{fig:collaboration-fairness-routes}(b)与(c)示意联合配送与转归解；具体分配数字见表\ref{tab:allocation}。

\subsection{动态需求响应层}

前述各节共有一个前提: 一日的订单在车辆出发之前已全部已知.
现实运行中并非如此——车辆已经发出, 新订单仍在陆续到达,
既有订单也可能被取消或修改需求量与时间要求.
此时已经走过的路径不能重来, 车辆的位置、载重、荷电状态,
以及已经发生的成本与排放同样不能清零.
本节把前述静态方案放入真实运行环境, 考察动态需求出现之后剩余路径应如何重新组织.

动态实验的三个对照臂为完全信息静态、机械在线和碳感知在线；完全信息静态预先知道全日订单，只作理论参照。现有动态结果包未通过验收且求解器没有进入，因此本节一个结果数字也不填写。正式结果产生时，三个臂均须同时报告完成客户数与完成需求量。

\begin{table}[htbp]
  \centering
  \caption{动态需求响应方式对比}
  \label{tab:dynamic}
  \scriptsize
  \resizebox{\textwidth}{!}{%
  \begin{tabular}{lrrrr}
    \toprule
    指标（单位） & 完全信息静态 & 机械在线 & 碳感知在线 & 碳感知相对机械在线\\
    \midrule
    单目标总账（元） & \TBD{动态正式批} & \TBD{动态正式批} & \TBD{动态正式批} & \TBD{比较}\\
    车辆固定成本（元） & \TBD{动态正式批} & \TBD{动态正式批} & \TBD{动态正式批} & \TBD{比较}\\
    非能源里程成本（元） & \TBD{动态正式批} & \TBD{动态正式批} & \TBD{动态正式批} & \TBD{比较}\\
    燃油成本（元） & \TBD{动态正式批} & \TBD{动态正式批} & \TBD{动态正式批} & \TBD{比较}\\
    电费（元） & \TBD{动态正式批} & \TBD{动态正式批} & \TBD{动态正式批} & \TBD{比较}\\
    碳价折算成本（元） & \TBD{动态正式批} & \TBD{动态正式批} & \TBD{动态正式批} & \TBD{比较}\\
    响应订单数（个） & \TBD{动态正式批} & \TBD{动态正式批} & \TBD{动态正式批} & \TBD{比较}\\
    拒绝/外包订单数（个） & \TBD{动态正式批} & \TBD{动态正式批} & \TBD{动态正式批} & \TBD{比较}\\
    实际车辆数（辆） & \TBD{动态正式批} & \TBD{动态正式批} & \TBD{动态正式批} & \TBD{比较}\\
    重规划次数（次） & \TBD{动态正式批} & \TBD{动态正式批} & \TBD{动态正式批} & \TBD{比较}\\
    总里程（km） & \TBD{动态正式批} & \TBD{动态正式批} & \TBD{动态正式批} & \TBD{比较}\\
    燃油车直接排放（kgCO$_2$e） & \TBD{动态正式批} & \TBD{动态正式批} & \TBD{动态正式批} & \TBD{比较}\\
    电动车充电排放（kgCO$_2$e） & \TBD{动态正式批} & \TBD{动态正式批} & \TBD{动态正式批} & \TBD{比较}\\
    完成客户数（个） & \TBD{动态正式批} & \TBD{动态正式批} & \TBD{动态正式批} & \TBD{比较}\\
    完成需求量（kg） & \TBD{动态正式批} & \TBD{动态正式批} & \TBD{动态正式批} & \TBD{比较}\\
    \bottomrule
  \end{tabular}
  }
\end{table}

\subsection{数值实验分析讨论}

分析上述各节结果, 可以按对两类主体的可操作性把本文的机制归为两组.

第一组是成本与排放同向的机制. 打开车场之间的客户归属边界后, 联合配送相对两家企业各自独立经营,
总成本由3922.079元降至3033.206元, 下降22.663\%, 总行驶里程下降45.5\%;
该项改善不要求企业追加车辆资产, 也不依赖外部价格信号,
其主要障碍是收益如何在参与方之间分配, 属于合作治理问题而非技术问题.
第二组是成本与排放反向的机制. 在既定配送排班下把充电时刻移向低碳时段,
可移动电量174.732~kWh, 充电环节排放下降59.082~kgCO$_2$e, 但电费增加26.525元;
在现行单位碳价0.075元/kgCO$_2$e下, 该项改变使企业的完整目标上升22.094元.
提高电动车派遣比例同样属于此组, 且因涉及车辆日固定成本, 其门槛\TBD{车型层门槛测算}.

由此可得两条针对性的启示. 对配送企业, 应优先实施第一组机制:
跨企业的客户责任重组与实体车多趟利用不增加资本开支, 且降本与减碳同向;
对第二组机制, 在现行价格信号下不宜期待企业自发实施, 其未实施并非管理疏失, 而是核算结果.
对政策制定者, 财政资源不必投向第一组, 而应投向第二组, 且需区分层次:
配对算式给出的充电时移盈亏平衡碳价为0.449元/kgCO$_2$e, 约为现行水平的六倍;
车型指派层的对齐碳价与之不同, 具体区间见表\ref{tab:carbon-price}与\TBD{车型层扫描结果}.
低于上述水平时, 提高碳价只增加企业负担而不改变其决策; 相关层次的减排潜力将持续闲置.

需要说明的是, 本文的时段电价与电网碳强度取自单一代表电网日,
客户责任归属为公开披露的构造情景, 结论的量级依赖于该情景;
但"价格信号与碳信号在时间上不同向, 因而充电时刻须作为独立决策进入路径优化"这一结构性结论,
不依赖于具体数值. 调整相应参数后, 本文模型可退化为单车场、单趟或纯燃油车队等经典VRP形式,
相应结论亦随之退化为已有研究的特例.

\section{结语}

本文研究时变电网碳强度下多家配送企业协作的动态混合车队路径优化问题,
考虑多车场协同、混合车队、实体车多趟、分时充电排放、成员结算参与和动态订单等因素,
建立可行配送趟—实体车排班两层模型,
并在成熟混合遗传搜索之上设计面向该问题的特化求解组件对其进行求解.
通过公开标准算例和中国自有算例检验算法的求解表现,
并沿混合车队、真实补能、时变碳强度、多车场协同、收益公平与动态运行的顺序组织数值实验,
逐层考察每一项现实因素是否改变原有的最优配送组织与路径选择.
结果表明: 跨企业重组客户责任使总成本下降22.663\%, 降本与减碳同向, 企业具备自发实施的动机;
而充电时刻与车型指派两层的降本与减碳方向相反,
在现行碳价下企业的最优选择是维持既有安排, 充电时移的盈亏平衡碳价为0.449元/kgCO$_2$e;
把静态方案放入动态运行环境后, 上述结论\TBD{动态臂对照结论}.

后续研究可进一步考察不同经营周期下的联盟收益分配,
并通过正式机制实验检验非线性充电曲线、动态订单处理节奏、充电排队和随机行驶时间的作用.

% ================================================================
% 参考文献
% ================================================================

\begin{thebibliography}{99}

\bibitem{ref:55} 国务院. 关于印发2030年前碳达峰行动方案的通知[EB/OL]. (2021-10-26)[2026-07-23].
\url{https://app.www.gov.cn/govdata/gov/202110/26/477466/article.html}.

\bibitem{ref:52} Li Z, Chen Z, Li H, Guan C, Zhong M. On the value of orderly electric vehicle charging in carbon emission reduction[J]. Transportation Research Part D: Transport and Environment, 2024, 135: 104383.

\bibitem{ref:53} Hu H, Qian B, Xiao Y, Tang J, Ou J, Lin X, He P, Zhang F. Low-carbon scheduling strategy for electric vehicles considering carbon emission flow and dynamic electricity prices[J]. Frontiers in Energy Research, 2024, 12: 1519963.

\bibitem{ref:54} Zhang J, Jing W, Lu Z, Wu H, Wen X. Collaborative strategy for electric vehicle charging scheduling and route planning[J]. IET Smart Grid, 2024, 7(5): 628--642.

\bibitem{ref:50} Cheng K W, Bian Y, Shi Y, Chen Y. Carbon-aware EV charging[C]//2022 IEEE International Conference on Communications, Control, and Computing Technologies for Smart Grids. Singapore: IEEE, 2022: 186--192.

\bibitem{ref:90} Miyabe R, Fujimoto Y, Hayashi Y. Low-carbon routing and charging planning for electric freight trucks utilizing local surplus solar power[J]. Journal of Energy Storage, 2025, 132: 117626.

\bibitem{ref:27} Goeke D, Schneider M. Routing a mixed fleet of electric and conventional vehicles[J]. European Journal of Operational Research, 2015, 245(1): 81--99.

\bibitem{ref:6} 李得成, 陈彦如, 张宗成. 基于分支定价算法的电动车与燃油车混合车辆路径问题研究[J]. 系统工程理论与实践, 2021, 41(4): 995--1009.\newline
Li D C, Chen Y R, Zhang Z C. A branch-and-price algorithm for electric vehicle routing problem with time windows and mixed fleet[J]. Systems Engineering --- Theory \& Practice, 2021, 41(4): 995--1009.

\bibitem{ref:23} 陈婉茹, 徐光明, 张得志, 等. 碳交易机制下多中心混合车队配送路径和速度优化研究[J]. 系统工程理论与实践, 2023, 43(11): 3320--3335.\newline
Chen W R, Xu G M, Zhang D Z, et al. Multi-depot mixed fleet routing and speed optimization under a carbon trading mechanism[J]. Systems Engineering --- Theory \& Practice, 2023, 43(11): 3320--3335.

\bibitem{ref:45} Keskin M, \c{C}atay B. Partial recharge strategies for the electric vehicle routing problem with time windows[J]. Transportation Research Part C: Emerging Technologies, 2016, 65: 111--127.

\bibitem{ref:46} Diefenbach H, Emde S, Glock C H. Multi-depot electric vehicle scheduling in in-plant production logistics considering non-linear charging models[J]. European Journal of Operational Research, 2023, 306(2): 828--848.

\bibitem{ref:73} Kullman N D, Froger A, Mendoza J E, Goodson J C. frvcpy: An open-source solver for the fixed route vehicle charging problem[J]. INFORMS Journal on Computing, 2021, 33(4): 1277--1283.

\bibitem{ref:7} Wang Y, Zhou J, Sun Y, Fan J, Wang Z, Wang H. Collaborative multidepot electric vehicle routing problem with time windows and shared charging stations[J]. Expert Systems with Applications, 2023, 219: 119654.

\bibitem{ref:25} Wang Y, Wei Z, Luo S, Zhou J, Zhen L. Collaboration and resource sharing in the multidepot time-dependent vehicle routing problem with time windows[J]. Transportation Research Part E: Logistics and Transportation Review, 2024, 192: 103798.

\bibitem{ref:49} Fern\'{a}ndez E, Roca-Riu M, Speranza M G. The shared customer collaboration vehicle routing problem[J]. European Journal of Operational Research, 2018, 265(3): 1078--1093.

\bibitem{ref:8} Soriano A, Gansterer M, Hartl R F. The multi-depot vehicle routing problem with profit fairness[J]. International Journal of Production Economics, 2023, 255: 108669.

\bibitem{ref:91} 饶卫振, 苗晓河, 朱庆华, 姜力文. 考虑企业服务质量差异的协作配送问题及成本分摊方法研究[J]. 系统工程理论与实践, 2022, 42(10): 2721--2739.

\bibitem{ref:93} Shi Y, Lin Y, Lim M K, Tseng M L. The bi-objective mixed-fleet vehicle routing problem under decentralized collaboration and time-of-use prices[J]. Expert Systems with Applications, 2025, 273: 126875.

\bibitem{ref:92} 饶卫振, 张云东, 刘从虎, 于灏, 侯艳辉. 一种求解协作配送成本分摊问题核仁解的近似迭代算法[J]. 系统工程理论与实践, 2019, 39(6): 1517--1534.

\bibitem{ref:9} Shi Y, Lin N, Han Q, Zhang T, Shen W. A method for transportation planning and profit sharing in collaborative multi-carrier vehicle routing[J]. Mathematics, 2020, 8(10): 1788.

\bibitem{ref:10} Agussurja L, Lau H C, Cheng S F. Achieving stable and fair profit allocation with minimum subsidy in collaborative logistics[C]//Proceedings of the AAAI Conference on Artificial Intelligence, 2016, 30(1): 3785--3792.

\bibitem{ref:14} Psaraftis H N. Dynamic vehicle routing problems[M]//Golden B L, Assad A A. Vehicle Routing: Methods and Studies. Amsterdam: North-Holland, 1988: 223--248.

\bibitem{ref:18} 李阳, 范厚明, 张晓楠. 动态需求下车辆路径问题的周期性优化模型及求解[J]. 中国管理科学, 2022, 30(8): 254--266.\newline
Li Y, Fan H M, Zhang X N. A periodic optimization model and solution for capacitated vehicle routing problem with dynamic requests[J]. Chinese Journal of Management Science, 2022, 30(8): 254--266.

\bibitem{ref:22} Gendreau M, Guertin F, Potvin J Y, et al. Parallel tabu search for real-time vehicle routing and dispatching[J]. Transportation Science, 1999, 33(4): 381--390.

\bibitem{ref:71} 邱莹莹. 低碳背景下考虑动态需求的混合车队物流配送路径问题研究[D]. 上海: 上海理工大学, 2024.\newline
Qiu Y Y. Research on the vehicle routing problem of mixed-fleet considering dynamic demand under the background of low carbon[D]. Shanghai: University of Shanghai for Science and Technology, 2024.

\bibitem{ref:44} Zhen L, Ma C, Wang K, Xiao L, Zhang W. Multi-depot multi-trip vehicle routing problem with time windows and release dates[J]. Transportation Research Part E: Logistics and Transportation Review, 2020, 135: 101866.

\bibitem{ref:51} 姜广田, 纪皎月, 董佳伟. 绿色物流配送下的多车型动态车辆路径优化[J]. 系统工程理论与实践, 2024, 44(7): 2362--2380.\newline
Jiang G T, Ji J Y, Dong J W. Multi-vehicle dynamic vehicle routing optimization in green logistics distribution[J]. Systems Engineering --- Theory \& Practice, 2024, 44(7): 2362--2380.

\bibitem{ref:48} 陈雨蝶, 干宏程, 程亮, 等. 双碳背景下复杂冷链物流模型及求解算法[J/OL]. 系统工程理论与实践. DOI: 10.12011/SETP2024-2027.\newline
Chen Y D, Gan H C, Cheng L, et al. Complex cold chain logistics model and its optimization algorithm under the background of double carbon[J/OL]. Systems Engineering --- Theory \& Practice. DOI: 10.12011/SETP2024-2027.

\bibitem{ref:36} Demir E, Bektas T, Laporte G. An adaptive large neighborhood search heuristic for the pollution-routing problem[J]. European Journal of Operational Research, 2012, 223(2): 346--359.

\bibitem{ref:47} Bektas T, Laporte G. The pollution-routing problem[J]. Transportation Research Part B: Methodological, 2011, 45(8): 1232--1250.

\bibitem{ref:12} Froger A, Jabali O, Mendoza J E, Laporte G. The electric vehicle routing problem with capacitated charging stations[J]. Transportation Science, 2022, 56(2): 460--482.

\bibitem{ref:13} 胡路, 乐诗彤, 朱娟秀. 考虑多充电桩排队和时间窗的电动货车路径规划[J]. 西南交通大学学报, 2025, 60(2): 299--307.\newline
Hu L, Le S T, Zhu J X. Electric truck route planning considering multiple charging pile queues and time windows[J]. Journal of Southwest Jiaotong University, 2025, 60(2): 299--307.

\bibitem{ref:76} Vidal T, Crainic T G, Gendreau M, Lahrichi N, Rei W. A hybrid genetic algorithm for multidepot and periodic vehicle routing problems[J]. Operations Research, 2012, 60(3): 611--624.

\bibitem{ref:77} Vidal T, Crainic T G, Gendreau M, Prins C. A unified solution framework for multi-attribute vehicle routing problems[J]. European Journal of Operational Research, 2014, 234(3): 658--673.

%% 已撤（正文不再引用路线池重组） \bibitem{ref:79} Rochat Y, Taillard \'E D. Probabilistic diversification and intensification in local search for vehicle routing[J]. Journal of Heuristics, 1995, 1(1): 147--167.

\bibitem{ref:80} Wang W, Adulyasak Y, Cordeau J F, He G. The heterogeneous-fleet electric vehicle routing problem with nonlinear charging functions[J]. Transportation Research Part C: Emerging Technologies, 2025, 170: 104932.

\bibitem{ref:64} Wouda N A, Lan L, Kool W. PyVRP: A high-performance VRP solver package[J]. INFORMS Journal on Computing, 2024, 36(4): 943--955.

\bibitem{ref:78} Vidal T, Crainic T G, Gendreau M, Prins C. Implicit depot assignments and rotations in vehicle routing heuristics[J]. European Journal of Operational Research, 2014, 237(1): 15--28.

\bibitem{ref:81} Huangfu Q, Hall J A J. Parallelizing the dual revised simplex method[J]. Mathematical Programming Computation, 2018, 10(1): 119--142.

\bibitem{ref:84} OpenStreetMap contributors. OpenStreetMap[DB/OL]. [2026-07-18]. \url{https://www.openstreetmap.org}.

\bibitem{ref:85} Luxen D, Vetter C. Real-time routing with OpenStreetMap data[C]//Proceedings of the 19th ACM SIGSPATIAL International Conference on Advances in Geographic Information Systems. New York: ACM, 2011: 513--516.

\bibitem{ref:87} 全国碳市场信息网. 全国碳排放权交易市场每日综合价格行情[DB/OL]. [2026-07-17]. \url{https://www.cets.org.cn/}.

\bibitem{ref:88} 广东省发展改革委. 关于我省新能源汽车用电价格有关问题的通知(粤发改价格〔2018〕313号)[EB/OL]. (2018-07-02)[2026-07-23].
\url{https://drc.gd.gov.cn/ywtz/content/post_833857.html}.

\bibitem{ref:83} 国家发展改革委. 关于第三监管周期省级电网输配电价及有关事项的通知(发改价格〔2023〕526号)[EB/OL]. (2023-05-15)[2026-07-23].
\url{https://www.ndrc.gov.cn/xxgk/zcfb/tz/202305/t20230515_1355747.html}.

\bibitem{ref:82} Li Y, Zhang S, Li W, Liu Y, Qin Y, Wei H, Kang C, Zhang N, Du E. High temporal and spatial resolution projected electricity carbon emission factors of China from 2025--2060[J]. Scientific Data, 2026, 13: 926. Dataset DOI: 10.6084/m9.figshare.28953545.

\bibitem{ref:74} Vidal T, Crainic T G, Gendreau M, Prins C. A hybrid genetic algorithm with adaptive diversity management for a large class of vehicle routing problems with time-windows[J]. Computers \& Operations Research, 2013, 40(1): 475--489.

\bibitem{ref:75} Lei Z, Hao J K. A GPU-accelerated hybrid method for a class of multi-depot vehicle routing problems[EB/OL]. arXiv:2605.05208, 2026.

\end{thebibliography}

\end{document}
